\chapter{外设、网络与人机接口}

\begin{keyidea}
介质和物理接口各不相同，但高吞吐外设普遍依赖队列、描述符、DMA 和完成通知。本章用网络与外设案例比较这些共同机制。
\end{keyidea}

\section{一、网卡与网络队列}

网卡面对的约束与磁盘相反：磁盘是“你求它快”，网络是“它不管你收
不收”。链路上的帧按线速到达，不能暂停、不可减速；对端何时发、
发多快，本机说了不算。网卡的全部队列设计都在回答同一个问题：帧
已经在线上了，CPU 还没准备好，怎么办。本节沿帧的生命周期走一
遍：线上怎么数帧，帧进内存走哪条环，CPU 以什么节奏被叫醒，多核
怎样分流，以及数据怎样少搬一次是一次。

\topic{帧、MTU 与线速}以太帧在线上的完整样子：先是 7 字节前导码
加 1 字节帧起始符（SFD），让接收方的时钟恢复电路锁定比特边界；
随后 6 字节目的 MAC、6 字节源 MAC、2 字节类型字段（EtherType），
指明载荷交给谁——IPv4、IPv6 还是 ARP；载荷 46--1500 字节，不足
46 字节必须填充；最后 4 字节 FCS，对整帧做 CRC32 校验，错帧在网
卡层面直接丢弃，不进内存。帧与帧之间还有 12 字节的帧间隙
（IFG），是链路上的静默间隔。MTU 1500 指的是载荷上限；线上一帧
实为 14 字节头加 1500 字节载荷加 4 字节 FCS，共 1518 字节，再算
上前导与帧间隙，每帧占线 1538 字节。

把这一帧的字节布局画出来，固定开销与载荷的比例就有了形状：

\begin{pdfdiagram}
  % 线上一帧：字段宽度为示意，总长对应 1538 字节
  \node[hw label, anchor=south west] at (0,3.3) {线上一帧共占 1538 字节（字段宽度为示意）};
  \draw[BookTeal, line width=0.8pt, fill=BookCodeBg] (0,2.3) rectangle (1.0,3.2);
  \node[align=center] at (0.5,2.75) {前导\\8 B};
  \draw[BookTeal, line width=0.8pt, fill=BookCodeBg] (1.0,2.3) rectangle (2.1,3.2);
  \node[align=center] at (1.55,2.75) {目的\\6 B};
  \draw[BookTeal, line width=0.8pt, fill=BookCodeBg] (2.1,2.3) rectangle (3.2,3.2);
  \node[align=center] at (2.65,2.75) {源\\6 B};
  \draw[BookTeal, line width=0.8pt, fill=BookCodeBg] (3.2,2.3) rectangle (4.1,3.2);
  \node[align=center] at (3.65,2.75) {类型\\2 B};
  \draw[BookAccent, line width=0.9pt, fill=BookCodeBg] (4.1,2.3) rectangle (9.0,3.2);
  \node at (6.55,2.75) {载荷 46--1500 B（MTU 1500）};
  \draw[BookTeal, line width=0.8pt, fill=BookCodeBg] (9.0,2.3) rectangle (9.9,3.2);
  \node[align=center] at (9.45,2.75) {FCS\\4 B};
  \draw[BookMuted, line width=0.8pt, fill=white] (9.9,2.3) rectangle (11.0,3.2);
  \node[align=center] at (10.45,2.75) {IFG\\12 B};
  % 开销占比对照：小帧 vs 大帧
  \node[hw label, anchor=east] at (-0.15,1.35) {最小帧 64 B};
  \draw[BookRed, line width=0.8pt, fill=white] (0,1.0) rectangle (2.0,1.7);
  \node[hw label] at (1.0,1.35) {开销 45\%};
  \draw[BookAccent, line width=0.8pt, fill=BookCodeBg] (2.0,1.0) rectangle (4.5,1.7);
  \node[hw label] at (3.25,1.35) {载荷};
  \node[hw label, anchor=west] at (4.7,1.35) {线上 84 B，载荷率约一半（548 Mb/s）};
  \node[hw label, anchor=east] at (-0.15,0.35) {最大帧 1518 B};
  \draw[BookRed, line width=0.8pt, fill=white] (0,0.0) rectangle (0.11,0.7);
  \draw[BookAccent, line width=0.8pt, fill=BookCodeBg] (0.11,0.0) rectangle (4.5,0.7);
  \node[hw label] at (2.3,0.35) {载荷};
  \node[hw label, anchor=west] at (4.7,0.35) {线上 1538 B，开销仅 2.5\%（975 Mb/s）};
\end{pdfdiagram}
\diagramnote{以太帧在线上的字节布局：前导与 SFD 8 字节用于时钟锁定，MAC 地址与类型字段 14 字节，载荷最多 1500 字节，FCS 4 字节校验，帧间还有 12 字节静默。每帧 38 字节的固定开销与帧长无关：64 字节小帧把它摊成 45\% 的负担，1518 字节大帧只付 2.5\%——压力测试用小帧考 pps、吞吐测试用大帧考带宽，原因就在这张图上。}

\begin{longtable}{L{0.20\linewidth}L{0.12\linewidth}L{0.56\linewidth}}
\toprule
字段 & 字节数 & 作用 \\
\midrule
前导码 + SFD & 8 & 时钟恢复与帧定界，属于物理层，不进内存 \\\tablerowrule
目的 MAC & 6 & 网卡按它做第一层过滤，不符则直接丢弃 \\\tablerowrule
源 MAC & 6 & 标识发送方，交换机按它学习转发表 \\\tablerowrule
类型 EtherType & 2 & 指明载荷协议：IPv4、IPv6、ARP 等 \\\tablerowrule
载荷 payload & 46--1500 & MTU 所指的部分，不足 46 字节要填充 \\\tablerowrule
FCS & 4 & CRC32 校验，错帧在网卡即丢弃 \\\tablerowrule
帧间隙 IFG & 12 & 帧间静默，链路占用但无内容 \\
\bottomrule
\end{longtable}

线速的换算从这里开始。1 GbE 的线上比特率是 $10^{9}$ bit/s——第九
章讲眼图时给过这个速率的直观：1 Gbps 下每个比特只有 1 ns 的判决窗
口。帧越小，每秒能过的帧越多，因为每帧都要付 38 字节的固定开销
（前导 8、帧间隙 12、头 14、FCS 4）。两个端点情形：最小帧 64 字
节，线上占 $64+8+12=84$ 字节即 672 bit，1 GbE 每秒过
$10^{9}/672\approx1.49\times10^{6}$ 帧，约 1.49 Mpps；最大帧 1518
字节，线上占 1538 字节即 12304 bit，每秒约
$10^{9}/12304\approx81274$ 帧。10 GbE 把这两个数都乘 10：约 14.88
Mpps 与约 81.3 万 pps。

\begin{longtable}{L{0.16\linewidth}L{0.20\linewidth}L{0.24\linewidth}L{0.24\linewidth}}
\toprule
帧长（字节） & 线上占用（字节） & 1 GbE（pps） & 10 GbE（pps） \\
\midrule
64（最小帧） & 84 & 约 149 万 & 约 1488 万 \\\tablerowrule
512 & 532 & 约 23.5 万 & 约 235 万 \\\tablerowrule
1518（最大帧） & 1538 & 81274 & 约 81.3 万 \\
\bottomrule
\end{longtable}

有效载荷效率值得单独算：小帧线速下，每秒载荷只有
$1.49\times10^{6}\times46\times8\approx548$ Mb/s，约线速的一半；
大帧下 $81274\times1500\times8\approx975$ Mb/s，接近满载。所以
“线速”与“pps”是两个独立的指标：线速（bit/s）是物理层能力，
pps 度量每帧固定成本的总量——网卡 DMA 一次、描述符消耗一个、协
议栈走一遍，都与帧长无关。压力测试用 64 字节小帧考 pps，吞吐测
试用 1518 字节大帧考带宽，两个数字各回答一个问题；选型时先量业
务的平均帧长，再决定该为哪个指标付钱。

帧格式也不是只有这一种长相。802.1Q 在源 MAC 与类型字段之间插入
4 字节标签，携带 VLAN 编号，帧长上限随之变成 1522 字节；QinQ、
MPLS 再往外套。每多一层封装，MTU 的有效载荷都会相应减少：隧道出
口看到的 1500 字节帧里，真正的数据可能只剩 1400 余字节——VPN
与隧道环境里“MTU 要调小”这条运维经验的来源，就是各层封装头的
总长度。

最小帧 64 字节也有来历。早期共享总线式以太网靠冲突检测
（CSMA/CD）工作：发送方必须发够长的时间，才能让冲突信号从总线
最远端传回来，否则冲突发生而发送方浑然不觉。按 10Mb/s 时代的线
缆长度与传播时延算，这个窗口定为 512 比特时，即 64 字节——最短
合法帧由此而来。交换式全双工以太网早已没有冲突，最小帧却作为约
定的一部分留了下来，今天我们数 pps 时仍在为它付费：每帧 38 字节
的固定开销，一半来自这段历史。

MTU 本身也可以调。数据中心内部常用 9000 字节载荷的巨型帧
（jumbo frame）：线上每帧约 9038 字节，1 GbE 每秒帧数降到约 1.4
万，每帧固定成本被摊到六倍的载荷上，pps 与中断、协议栈开销同步
下降约六倍。前提是路径上每一跳的 MTU 都得一致放大，否则分片甚
至静默丢包会抵消全部收益——巨型帧是全网约定，不是单机参
数。

“网卡通常不能要求发送端立即停止”严格说有一个例外：以太网流控允许接收方发 PAUSE
帧，请对端暂停发送指定时长。它的本意是防丢包，实际效果是把缓冲
区不足的问题传染给上游——一台主机收不动，整条链路上的无辜流量
一起被压。现代数据中心普遍关闭全局 PAUSE，丢包问题留给更深的
环、更快的消费和拥塞协议去解决；存储网络使用的按类暂停（PFC）
是另一套权衡，以队头阻塞为代价换无损。这个例外恰好坐实本节的出
发点：默认情况下，到达速率不由接收方控制，缓冲深度与消费速率必
须自行处理恢复。

\topic{接收环与发送环}网卡与内存之间走的就是存储与设备事务相关章节的描述符环，
机制原封不动：环是内存里的描述符数组，CPU 推进软件指针、设备推
进硬件指针，doorbell 只提示“指针动了”，完成以设备写回的状态位
为准。网卡的特殊性在于，收发两个方向的环各有一条不对称的约定：
接收环要预投，发送环要回收。

RX 环是预投的。收帧不可预约——对端随时发，网卡通常不能要求发送端立即停止——所
以驱动必须提前把一批空缓冲区的地址铺满接收环。帧到达后，网卡
MAC 校验 FCS，由 DMA 把帧直接写进下一个空缓冲区，写回描述符状
态（帧长、校验结果、RSS 队列号），再按中断策略通知 CPU。驱动取
走数据、把缓冲区交给协议栈之后，必须立刻补投新的空缓冲区：预投
是义务，不是优化。环上可用缓冲区一旦耗尽，再到的帧没有地方放，
只能丢弃——帧本身没有错，是没人接。

TX 环是回收的。发送由软件发起：驱动把待发帧的描述符挂上发送
环、敲 doorbell；网卡按描述符 gather 数据、送上线路，发完写回完
成状态；驱动在中断或轮询里回收描述符并释放缓冲区。回收不及时，
环被挂满，新帧无处提交，协议栈只能排队等待，背压沿协议栈向上传
导。两个方向的环因此有不同的“满”的含义：RX 环空是丢包，TX 环
满是阻塞。

发送环的回收节奏同样可调：每发完一帧就中断一次太奢侈，驱动通常
让网卡攒一批完成再通知，或干脆在下次发送与定时轮询时顺手回收。
回收延迟直接占用环深——一批未回收的描述符，与一串未补投的空缓
冲区，是同一个问题的两面：环的有效深度，等于标称深度减去滞留在
“已完成未回收”状态里的那部分。

环深与丢包的关系可以直接算。RX 环的深度必须吸收“最坏通知延
迟”内到达的全部帧：1 GbE 小帧每秒约 149 万帧，若中断聚合加调度
把最坏通知延迟拖到 200 $\mu$s，期间到达
$1.49\times10^{6}\times200\times10^{-6}\approx298$ 帧——环深 64
必然丢，512 才够；10 GbE 同样的窗口要容纳约 2980 帧，常见取
4096。丢包在以太网层面不报错、不重传，只体现为网卡计数器里的数
字，最后由 TCP 重传表现为“网速偶尔抖一下”。存储与设备事务相关章节“缓冲区何时
归 CPU、何时归设备”的所有权问题，在这里直接变成丢包与否的分界
线。

\begin{lstlisting}
$ ethtool -g eth0      # show ring depths: RX 512, TX 512
$ ethtool -S eth0 | grep -E 'rx_missed|rx_no_buffer'
     rx_missed_errors: 1382     # frames arrived with no buffer
$ ethtool -G eth0 rx 4096      # deepen RX ring (driver permitting)
\end{lstlisting}

丢包计数是这条链上最诚实的仪表。应用看到的是吞吐波动与重传，协
议栈看到的是乱序，只有网卡计数器直接说出“哪一环没接住”：
rx\_no\_buffer 指预投断供，rx\_missed 指环深或处理速度不够。排障
的顺序因此固定：先看计数器分辨丢在哪一层，再回头调环深、中断参
数或核数——本节后面三个话题，正好各管其中一环。

\topic{中断聚合}每帧一中断在小帧线速下根本不成立。1.49 Mpps 意味
着每 0.67 $\mu$s 一次中断；一次完整的中断处理——保存现场、读状态、
递交协议栈、恢复现场——按保守的 3 $\mu$s 计，单核每秒最多服务约 33
万次。需要的算力是 $1.49\times10^{6}\times3\mu s\approx4.5$ 颗核
的全部时间，而且全花在中断进出上，协议栈和应用一行代码都轮不
到：越忙越处理不完，吞吐反而崩，这就是中断风暴（receive
livelock）。0.67 $\mu$s 这个帧间隔值得记住：它就是“软件为每帧固定成
本付费”的时钟节拍。

把帧间隔与处理时长画在同一条时间轴上，风暴的自我放大就
清楚了：

\begin{waveform}[x=0.8cm,y=0.9cm]
  % 帧到达：每 0.67 µs 一帧
  \foreach \x in {1,2,...,12} \draw[BookTeal, line width=0.9pt] (\x,5.9) -- (\x,6.5);
  \node[hw label, anchor=east] at (0.8,6.2) {帧到达};
  % 3 µs 里又到约 4 帧
  \draw[BookMuted, line width=0.6pt, <->] (1,5.3) -- (5.5,5.3);
  \node[hw label, anchor=south] at (3.25,5.4) {3 $\mu$s 里又到约 4 帧};
  % CPU：一次中断处理 3 µs，处理一帧又积下数帧
  \draw[BookAccent, fill=BookCodeBg, line width=0.8pt] (1,4.3) rectangle (5.5,4.9);
  \node at (3.25,4.6) {处理第 1 帧：3 $\mu$s};
  \draw[BookAccent, fill=BookCodeBg, line width=0.8pt] (5.5,4.3) rectangle (10,4.9);
  \node at (7.75,4.6) {处理第 2 帧：3 $\mu$s};
  \draw[BookAccent, fill=BookCodeBg, line width=0.8pt] (10,4.3) rectangle (13,4.9);
  \node[hw label] at (11.5,4.6) {处理第 3 帧 $\cdots$};
  \node[hw label, anchor=east] at (0.8,4.6) {中断处理};
  % 未处理积压：阶梯单调上涨（每个 3 µs 窗口净增约 3 帧）
  \draw[BookRed, line width=0.9pt] (1,1.2) -- (1,1.47) -- (2,1.47) -- (2,1.74) -- (3,1.74) -- (3,2.01) -- (4,2.01) -- (4,2.28) -- (5,2.28) -- (5,2.55) -- (5.5,2.55) -- (5.5,2.28) -- (6,2.28) -- (6,2.55) -- (7,2.55) -- (7,2.82) -- (8,2.82) -- (8,3.09) -- (9,3.09) -- (9,3.36) -- (10,3.36) -- (10,3.09) -- (11,3.09) -- (11,3.36) -- (12,3.36) -- (12,3.63) -- (13,3.63);
  \node[hw label, anchor=east] at (0.8,2.4) {未处理积压};
  \node[hw label] at (6.6,0.55) {帧每 0.67 $\mu$s 一帧，处理一帧又净增约 3 帧：越忙越处理不完};
\end{waveform}
\diagramnote{中断风暴（receive livelock）的时序（1 GbE 小帧线速）。帧每 0.67 $\mu$s 到达一帧，一次完整的中断处理——保存现场、读状态、递交协议栈、恢复现场——按保守的 3 $\mu$s 计：处理一帧的功夫线上又到约四帧，未处理积压单调上涨，需要的算力约 $1.49\times10^{6}\times3\ \mu$s 即 4.5 颗核的全部时间，协议栈和应用一行代码都轮不到。越忙越处理不完、吞吐反而崩；出路是把到达率除以批大小，即下一段的中断聚合。}

NAPI 把“中断驱动”改成“中断提示加轮询收尾”：第一帧触发中断后
关闭该队列的中断，驱动进入轮询，在预算内尽量多地收完积压的帧，
收干了再重新开中断——忙时纯轮询，闲时纯中断，两种状态自动切
换。网卡硬件侧的中断合并（interrupt moderation）提供两个阈值：
攒够 N 帧发一次中断，或距上次中断满 T $\mu$s 发一次，先到为准；低负
载时靠超时保证单帧不被无限积压。

补一个实现细节：NAPI 的轮询收包运行在软中断（softirq）上下文，
协议栈的 IP、TCP 处理大多也在其中。top 里看到的 si 一项的占比，
就是这条收包路径的 CPU 开销；它高到挤压应用，说明帧处理本身—
—而不只是中断进出——已经成为瓶颈，下一步手段是多队列分流，见
下一段。

权衡可以量化。设合并参数 N=32、T=40 $\mu$s：小帧线速下 32 帧只需
$32\times0.672\approx21.5\mu s$，先于超时触发，中断率从约 149 万
次/秒降到 $1.49\times10^{6}/32\approx4.7$ 万次/秒，CPU 开销
$4.7\times10^{4}\times3\mu s\approx0.14$ 颗核——从中断风暴的 4.5
颗核降到约七分之一颗。代价是延迟：一帧平均多等约半批的时间，十
微秒量级。吞吐型业务把 N、T 调大；低延迟业务（行情、内存数据
库）调小甚至关闭合并，改用绑核加 busy-poll——拿一颗核 100\% 的
占用，换确定的几微秒时延。两个方向都对，区别只在业务把哪一头算
作成本。

把两种策略各画一条时序，这笔交换的形状如下：

\begin{waveform}[x=0.8cm,y=0.9cm]
  % 参考虚线：第 32 帧到达触发 / 等满 T 超时触发
  \draw[BookRule, line width=0.45pt, dashed] (8.75,1.35) -- (8.75,3.9);
  \draw[BookRule, line width=0.45pt, dashed] (12.0,1.35) -- (12.0,3.9);
  % 帧到达（线速，不聚合场景）
  \foreach \x in {1.0,1.25,...,8.9} \draw[BookTeal, line width=0.9pt] (\x,7.0) -- (\x,7.55);
  \node[hw label, anchor=east] at (0.8,7.28) {帧到达（线速）};
  % 中断·不聚合：每帧一次
  \draw[BookAccent, line width=0.9pt] (1.0,5.4) -- (9.0,5.4);
  \foreach \x in {1.0,1.25,...,8.9} \draw[BookAccent, line width=0.9pt] (\x,5.4) -- (\x,6.0) -- ({\x+0.13},6.0) -- ({\x+0.13},5.4);
  \node[hw label, anchor=east] at (0.8,5.7) {中断·不聚合};
  \node[hw label] at (6.4,4.9) {不聚合：每帧一次，$\approx$149 万次/秒，CPU $\approx$4.5 核};
  % 帧到达（聚合场景：线速 32 帧 + 低负载单帧）
  \foreach \x in {1.0,1.25,...,8.9} \draw[BookTeal, line width=0.9pt] (\x,3.2) -- (\x,3.75);
  \draw[BookTeal, line width=0.9pt] (10.2,3.2) -- (10.2,3.75);
  \node[hw label, anchor=east] at (0.8,3.48) {帧到达};
  \draw[BookMuted, line width=0.6pt, <->] (1.0,4.05) -- (8.75,4.05);
  \node[hw label] at (4.9,4.34) {攒满 32 帧 $\approx$21.5 $\mu$s（先于超时）};
  \draw[BookMuted, line width=0.6pt, <->] (10.2,4.05) -- (12.0,4.05);
  \node[hw label] at (11.1,4.34) {等满 $T$=40 $\mu$s};
  \node[hw label] at (10.4,2.92) {低负载：单帧到达};
  % 中断·聚合：攒满一批或超时，才发一次
  \draw[BookAccent, line width=0.9pt] (1.0,1.6) -- (8.75,1.6) -- (8.75,2.2) -- (9.05,2.2) -- (9.05,1.6) -- (12.0,1.6) -- (12.0,2.2) -- (12.3,2.2) -- (12.3,1.6) -- (12.6,1.6);
  \node[hw label, anchor=east] at (0.8,1.9) {中断·聚合};
  \node[hw label] at (6.3,0.95) {聚合 $N$=32 / $T$=40 $\mu$s：$\approx$4.7 万次/秒，CPU $\approx$0.14 核};
\end{waveform}
\diagramnote{中断聚合的时序对比。上方不聚合：每帧一次中断，1.49 Mpps 折合每秒约 149 万次中断进出，约 4.5 颗核的算力全花在保存与恢复现场上。下方聚合：攒满 32 帧（线速下约 21.5 $\mu$s，先于超时）或等满 T=40 $\mu$s——先到为准——才发一次中断，中断率降到约 4.7 万次/秒，CPU 开销约 0.14 核；代价是一帧平均多等约半个聚合窗口。}

\begin{longtable}{L{0.16\linewidth}L{0.32\linewidth}L{0.20\linewidth}L{0.20\linewidth}}
\toprule
参数 & 机制 & 偏向吞吐 & 偏向时延 \\
\midrule
帧数阈值 N & 攒够 N 帧发一次中断 & 调大（如 32--64） & 调小或置 1 \\\tablerowrule
超时 T & 距上次中断满 T $\mu$s 即发 & 调大（几十 $\mu$s） & 调小或置 0 \\\tablerowrule
轮询预算 & NAPI 每次最多收多少帧 & 调大 & 适中，防占核过久 \\\tablerowrule
busy-poll & 应用自旋收帧，跳过中断 & 浪费核 & 时延最低且稳 \\
\bottomrule
\end{longtable}

观察手段是现成的：每秒读一次 /proc/interrupts 里各网卡队列的中
断计数，与流量对比就知道合并参数是否合身——中断率贴近 pps，说
明合并未生效；中断率过低而时延投诉上升，说明 N、T 过了头。
ethtool -c 与 -C 查询、调整这些参数。调网卡的第一步通常不是换硬
件，而是把中断率、环深、pps 三个数摆在一起看；三者相符，瓶颈才
轮到协议栈和应用。

\begin{lstlisting}
$ grep eth0 /proc/interrupts
 28:  1204431        0        0        0  IO-APIC  28-fasteoi  eth0-rx-0
 29:        0  1188204        0        0  IO-APIC  29-fasteoi  eth0-rx-1
 30:        0        0  1210022        0  IO-APIC  30-fasteoi  eth0-rx-2
 31:        0        0        0  1195330  IO-APIC  31-fasteoi  eth0-rx-3
# per-queue counters spread over 4 CPUs: RSS and affinity agree
\end{lstlisting}

读法：每个队列的中断计数应大致均匀地散在各核上；挤在一列说明
RSS 或绑核没有配对，整行全为 0 的队列说明中断亲和漏配。这个文
件也是下一话题的入口：多队列不是插上就有，是用对才有。

\topic{多队列与 RSS}单队列的天花板先算清楚：一个环只能由一颗核
服务——保序与 cache 局部性都要求如此——单核跑协议栈每秒约
1--2M 帧已是优化后的水平；10 GbE 小帧是 14.88 Mpps，差一个数量
级。单队列网卡插上万兆链路，瓶颈不在线，在环。用“冒险、异常与性能”一章 Little
定律的读法：收帧路径的吞吐等于并行消费者数乘单消费者速率，在途
帧排得再多，消费者只有一颗核。

多队列网卡把 RX/TX 环做成很多对，常见 16--128 对，每对绑一个
MSI-X 中断向量、绑一颗核。接收侧的分流由硬件完成：RSS 对每帧的
五元组（源 IP、目的 IP、源端口、目的端口、协议号）计算 Toeplitz
哈希，取哈希低位查间接表，把帧定向到对应队列。同一条流哈希相
同，永远进同一队列，天然保序；不同的流散到不同核，并行处理。发
送侧对称：软件按流或按 CPU 选发送队列，各核各挂各的环，提交路
径无锁。

分流的全过程画成一张图，就是哈希定向加按核绑环：

\begin{pdfdiagram}
  \node[hw label, anchor=south] at (5.5,0.68) {同一条流哈希相同，永远进同一队列——天然保序};
  \node[hw state, minimum width=2.4cm] (nic) at (1.5,-1.28) {网卡收包\\10GbE 小帧\\14.88 Mpps};
  \node[hw compute, minimum width=2.8cm] (hash) at (4.35,-1.28) {RSS：五元组\\Toeplitz 哈希\\低位查间接表};
  \draw[hw bus] (nic.east) -- (hash.west);
  % 队列（示意 4 行，常见 16--128 对）
  \node[hw memory, minimum width=1.6cm, minimum height=0.55cm, inner sep=1pt] (q0) at (7.05,0) {RX 环 0};
  \node[hw memory, minimum width=1.6cm, minimum height=0.55cm, inner sep=1pt] (q1) at (7.05,-0.85) {RX 环 1};
  \node[hw label] at (7.05,-1.7) {$\vdots$};
  \node[hw memory, minimum width=1.6cm, minimum height=0.55cm, inner sep=1pt] (q15) at (7.05,-2.55) {RX 环 15};
  % 核
  \node[hw compute, minimum width=1.3cm, minimum height=0.55cm, inner sep=1pt] (c0) at (9.6,0) {核 0};
  \node[hw compute, minimum width=1.3cm, minimum height=0.55cm, inner sep=1pt] (c1) at (9.6,-0.85) {核 1};
  \node[hw label] at (9.6,-1.7) {$\vdots$};
  \node[hw compute, minimum width=1.3cm, minimum height=0.55cm, inner sep=1pt] (c15) at (9.6,-2.55) {核 15};
  % 哈希到队列：共享竖直干线，分层折出
  \draw[hw bus] (hash.east) -- (6.05,-1.28);
  \draw[BookMuted, line width=0.62pt] (6.05,0) -- (6.05,-2.55);
  \draw[hw bus] (6.05,0) -- (q0.west);
  \draw[hw bus] (6.05,-0.85) -- (q1.west);
  \draw[hw bus] (6.05,-2.55) -- (q15.west);
  % 队列到核：一对一
  \draw[hw bus] (q0.east) -- (c0.west);
  \draw[hw bus] (q1.east) -- (c1.west);
  \draw[hw bus] (q15.east) -- (c15.west);
  \node[hw label] at (5.5,-3.4) {每对环绑一个 MSI-X 向量、绑一颗核；单核约 1.5 Mpps，16 颗核并行才接得住 14.88 Mpps};
\end{pdfdiagram}
\diagramnote{RSS 接收分流：网卡对每帧的五元组算 Toeplitz 哈希，取低位查间接表，把帧定向到对应 RX 环；每对环绑一个 MSI-X 中断向量、绑一颗核，各核独立处理自己环上的帧，提交路径无锁。同一条流哈希相同、永远进同一队列，保序由哈希天然保证；并行度由队列对数决定——10 GbE 小帧 14.88 Mpps 对单核约 1.5 Mpps，16 队列绑 16 核才接得住。}

配置按一遍数：10 GbE 小帧 14.88 Mpps，单核约 1.5 Mpps，需要约 10
颗核起步，留余量配 16 队列、16 个 MSI-X 向量、绑 16 颗核。三个
配置必须一起看：队列数、间接表、中断亲和——队列开了 32 条而中
断全绑在核 0，等于单队列；irqbalance 之类的自动均衡在万兆以上
反而可能降低性能，绑核要按队列手工固定。

同样的算法推到更高速率：25 GbE 小帧约 37.2 Mpps，40 GbE 约 59.5
Mpps，100 GbE 约 148.8 Mpps——单核约 1.5 Mpps 的处理率面前，需
要的核数按比例涨到约 25、40、100 颗，已经没有机器装得下。速率
每翻一番，“每帧固定成本”问题就尖锐一分：100 GbE 时代还按“中断
加协议栈逐帧处理”的模型工作是不成立的，内核旁路与硬件卸载从可
选项变成必选项，这正是零拷贝一段要收尾的方向。

\begin{longtable}{L{0.18\linewidth}L{0.36\linewidth}L{0.34\linewidth}}
\toprule
组件 & 作用 & 配置不当的典型症状 \\
\midrule
队列对 & 收发各一条环，独立 DMA 与中断 & 队列少于核数，部分核闲置 \\\tablerowrule
MSI-X 向量 & 每队列一个中断号 & 向量不足，多队列共享中断互拖 \\\tablerowrule
RSS 间接表 & 哈希值到队列的映射 & 映射不均，各核负载歪斜 \\\tablerowrule
中断亲和 & 把向量绑到指定核 & 全绑核 0，等于退回单队列 \\
\bottomrule
\end{longtable}

RSS 也有边界。它只认五元组：GRE、VxLAN 等隧道的外层包头都相
同，所有隧道流量挤进一条队列，较新的网卡支持按内层字段哈希来解
决。单条单条高带宽流永远只能跑一颗核的速度，这是保序的代价，只能靠多
开几条流来绕。与 RSS 同族的卸载还有 TSO 与 GRO：发送侧由网卡把
大块数据切成 MTU 帧，接收侧把同流相邻帧拼成大块再上交——它们
不改队列结构，改的是“每帧固定成本”里的帧数，与巨型帧是同一思
想的软硬件两种实现。

\topic{零拷贝的基本思想}先算拷贝一次的成本。传统收包路径上，网
卡 DMA 把帧写进内核缓冲区，这一步不占 CPU；协议栈处理后再由
copy\_to\_user 把数据复制一份到应用缓冲区——这一次是 CPU 逐字
节的复制，读一遍、写一遍。10 Gb/s 的流量约 1.25 GB/s，一次拷贝产
生 2.5 GB/s 的内存流量；收发各拷一次就是 5 GB/s，约占一条
DDR4-3200 通道（25.6 GB/s）可用带宽的两成，还没算 cache 被冲刷的
间接代价。用“冒险、异常与性能”一章 roofline 的话说：每次无谓的拷贝都在消耗全机
共享的带宽屋顶，留给应用的屋顶相应变矮。

零拷贝的思路是让数据在内存里只躺一份，流转的是“这块缓冲区归
谁”。sendfile 让文件页缓存直接对接 socket，配合支持
scatter-gather 的网卡，连“协议头加数据”的拼接都由网卡 gather
DMA 完成，CPU 一个字节不搬；splice 在两个文件描述符之间搬运页
指针；mmap 让应用直接读写内核页，省掉复制但引入页错误与同步的
成本；io\_uring 的固定缓冲区与 AF\_XDP 更进一步，把缓冲区的所有
权约定直接开放给用户态。注意“零拷贝”省的从来不是 DMA 那一次
——设备搬运本来就绕不开——省的是内核与用户态之间那一次 CPU
复制。

\begin{lstlisting}
#include <sys/sendfile.h>

/* stream a file to a socket without copying through user space */
ssize_t sendfile(int out_fd, int in_fd, off_t *offset, size_t count);

/* kernel path: page cache pages -> socket buffer -> NIC gather DMA;
   file contents never pass through a CPU register */
\end{lstlisting}

两条路径并排画出来，零拷贝省掉的是哪一段就一目了然：

\begin{pdfdiagram}
  % 传统路径：两次 CPU 拷贝
  \node[hw label, anchor=west] at (0.4,0.98) {传统：页缓存 $\to$ 应用缓冲 $\to$ socket 缓冲，两次 CPU 拷贝};
  \node[hw memory, minimum width=2.0cm, minimum height=0.8cm] (pc1) at (1.4,0) {页缓存};
  \node[hw memory, minimum width=2.0cm, minimum height=0.8cm] (ub) at (4.2,0) {应用缓冲};
  \node[hw memory, minimum width=2.0cm, minimum height=0.8cm] (sb1) at (7.0,0) {socket 缓冲};
  \node[hw memory, minimum width=2.0cm, minimum height=0.8cm] (nic1) at (9.8,0) {网卡};
  \draw[hw bus] (pc1.east) -- (ub.west);
  \draw[hw bus] (ub.east) -- (sb1.west);
  \draw[hw bus] (sb1.east) -- (nic1.west);
  \node[hw label, anchor=south] at (2.8,0.62) {① CPU 拷贝};
  \node[hw label, anchor=south] at (5.6,0.62) {② CPU 拷贝};
  \node[hw label, anchor=south] at (8.4,0.62) {③ DMA};
  % 零拷贝路径
  \node[hw label, anchor=west] at (0.4,-0.98) {零拷贝 sendfile：流转的是缓冲区归属，CPU 不搬数据};
  \node[hw memory, minimum width=2.0cm, minimum height=0.8cm] (pc2) at (1.4,-1.95) {页缓存};
  \node[hw memory, minimum width=2.6cm, minimum height=0.8cm] (sb2) at (5.6,-1.95) {socket 缓冲\\只记指针};
  \node[hw memory, minimum width=2.0cm, minimum height=0.8cm] (nic2) at (9.8,-1.95) {网卡};
  \draw[hw return] (pc2.east) -- (sb2.west);
  \draw[hw bus] (sb2.east) -- (nic2.west);
  \node[hw label, anchor=south] at (2.85,-1.33) {只改归属};
  \node[hw label, anchor=south] at (8.35,-1.33) {gather DMA};
  \node[hw label] at (5.5,-2.9) {10 Gb/s 下两次 CPU 拷贝产生 5 GB/s 内存流量，约占一条 DDR4-3200 通道（25.6 GB/s）的两成};
\end{pdfdiagram}
\diagramnote{发送一个文件的两种路径。传统路径上数据两次经过 CPU：页缓存拷进应用缓冲、应用缓冲再拷进 socket 缓冲，每次都是读一遍写一遍，10 Gb/s 的流量就此产生 5 GB/s 的内存流量。sendfile 让页缓存直接对接 socket，缓冲区只改归属、内容原地不动，网卡用 gather DMA 把协议头与数据拼接发出，CPU 一个字节不搬——省的是内核与用户态之间那次复制，设备 DMA 那一步本来就有。}

\begin{longtable}{L{0.18\linewidth}L{0.34\linewidth}L{0.34\linewidth}}
\toprule
机制 & 省掉的是哪次拷贝 & 适用场景 \\
\midrule
sendfile & 文件页缓存到用户态的往返 & 静态文件服务、CDN \\\tablerowrule
splice & 任意两个 fd 之间的搬运 & 代理、转发管道 \\\tablerowrule
mmap & 读方向的复制 & 随机访问大文件 \\\tablerowrule
io\_uring 固定缓冲 & 每次 I/O 的缓冲区注册与复制 & 高频异步 I/O \\\tablerowrule
AF\_XDP / DPDK & 整个内核协议栈路径 & 极端 pps 的专用应用 \\
\bottomrule
\end{longtable}

工程含义与边界。零拷贝的本质是所有权转移代替内容复制，正是第五
章“缓冲区何时归 CPU、何时归设备”那条问题在协议栈内部的延长
线。该机制也会引入额外代价：页共享带来生命周期管理与写时复制的复杂度，
数据量小时固定开销可能盖过节省；它适用的条件是大块数据、高带
宽、CPU 比带宽更紧张——这正是存储服务器、CDN、反向代理人手一
份 sendfile 的原因。从本章视角看，零拷贝与中断聚合、RSS 是同一
族手段：网络栈的性能工作，大半是在给“每帧固定成本”与“每字节
搬运成本”这两项开销找出路。

旁路也不是没有代价。把协议栈整个绕开，意味着 TCP 的拥塞控制、
分片、校验和、防火墙一并让位，应用要么自己实现，要么声明不需
要；内核协议栈占的 CPU，换来的正是这些不免费的功能。工程上成熟
的做法是分层：极端 pps 的入口（行情网关、负载均衡）用旁路或
AF\_XDP，业务主体仍走内核栈加零拷贝。存储与设备事务相关章节那句提醒在这里同样
成立：约定一旦绕开，原来由内核代管的所有权、可见性与错误处理，
全部换成应用自己的功课。

\section{二、外设队列综合案例与尾延迟}

前四节分别看了三类设备的介质与各自的队列结构，本节把它们汇合到一个
问题上：队列机制怎样决定外设实际表现出来的性能。

\topic{NVMe 的队列模型}
存储与设备事务相关章节讲描述符环时给过一般结构：环不动、指针动，doorbell 只通知、不
携带数据，完成与否以设备写回的状态为准。NVMe 把这套结构定型为规范，
并针对多核做了两处关键改动。第一处改动是队列成对：NVMe 的基本单元
是一对队列——提交队列 SQ（submission queue）与完成队列 CQ
（completion queue）。SQ 项是 64 字节的命令：操作码（读、写、
flush）、起始 LBA、块数，以及 PRP 数据指针列表，逐页给出数据在主机
内存里的物理地址；CQ 项是 16 字节的完成记录：命令 ID、状态码，以及
一个相位位——队列回绕一圈相位翻一次，主机据此区分新旧完成项，不必
比较指针。提交与完成分成两个环，各走各的指针：设备消费 SQ 不必等主
机回收 SQ，完成信息也不挤在命令项里写回。这是比存储与设备事务相关章节通用环更彻底
的分工。

一次提交的完整序列有七步，逐步对应存储与设备事务相关章节的每一条约定：

\begin{lstlisting}
// 注释（推进）：逐行追踪发起者、所有权和可见性；请求被接受不等于事务完成。
// 注释（边界）：以采样沿、状态位或 completion 为准，再允许消费者使用结果。
1. 驱动把 64 B 命令项写入本核 SQ 的尾槽
2. 屏障：保证命令项先于通知对设备可见
3. 写 SQ 尾 doorbell（一次 MMIO 写，内容只是新尾指针）
4. 设备看到指针差，DMA 取走新命令项
5. 设备执行：查 FTL、读 NAND，按 PRP 地址把数据 DMA 进内存
6. 设备写 16 B CQ 完成项，相位位翻转
7. 设备发 MSI-X 中断；驱动读 CQ、按命令 ID 找回请求并回收，
   最后写 CQ 头 doorbell，把完成槽位还给设备
\end{lstlisting}

第 3 步的 doorbell 仍然不携带数据：设备靠首尾指针的差补读新命令，偶
尔错过一次 doorbell 并无妨碍；第 6、7 步的次序也不许交换：先写完成
项、再发中断，主机在中断处理里读到的一定是已经写好的 CQ 项。存储与设备事务相关章节
“描述符先可见、再通知”和“完成以写回状态为准”这两条规则，在这里原
样成立，只是被写进了设备规范，由硬件而不是驱动员的细心来保证。

把这条链路画成图，主机内存、doorbell 寄存器与设备之间的
分工如下：

\begin{pdfdiagram}
  \node[hw state, minimum width=1.9cm] (cpu) at (0.9,0.55) {CPU\\（本核驱动）};
  % 主机内存：SQ 环与 CQ 环
  \draw[BookAmber, line width=0.42pt, rounded corners=2pt] (2.6,-1.3) rectangle (7.2,1.35);
  \node[hw label, anchor=west] at (2.75,1.08) {SQ：64 B 命令项环};
  \node[hw label, anchor=east] at (7.1,0.12) {主机内存（DRAM）};
  \node[hw state, minimum width=0.78cm, minimum height=0.5cm] (sq1) at (3.35,0.55) {命令};
  \node[hw state, minimum width=0.78cm, minimum height=0.5cm] at (4.35,0.55) {命令};
  \node[hw state, minimum width=0.78cm, minimum height=0.5cm] at (5.35,0.55) {命令};
  \node[hw state, minimum width=0.78cm, minimum height=0.5cm] (sq4) at (6.35,0.55) {$\cdots$};
  \node[hw state, minimum width=0.78cm, minimum height=0.5cm] at (3.35,-0.55) {完成};
  \node[hw state, minimum width=0.78cm, minimum height=0.5cm] at (4.35,-0.55) {完成};
  \node[hw state, minimum width=0.78cm, minimum height=0.5cm] at (5.35,-0.55) {完成};
  \node[hw state, minimum width=0.78cm, minimum height=0.5cm] (cq4) at (6.35,-0.55) {$\cdots$};
  \node[hw label, anchor=west] at (2.75,-1.05) {CQ：16 B 完成项环};
  \node[hw control, minimum width=1.6cm] (db) at (8.3,0.45) {doorbell\\寄存器};
  \node[hw compute, minimum width=1.9cm] (dev) at (10.45,0.45) {NVMe 控制器\\FTL + DMA};
  % ① 驱动把命令项写进本核 SQ 尾槽
  \draw[hw bus] (cpu.east) -- (sq1.west);
  \node[hw label, anchor=west] at (2.75,0.12) {① 写命令项};
  % ② 屏障后写 SQ 尾 doorbell（MMIO，控制）
  \draw[hw return] (cpu.south) -- (0.9,-2.0) -- (8.3,-2.0) -- (db.south);
  \node[hw label, above] at (4.5,-2.0) {② 写 SQ 尾 doorbell（MMIO）};
  % ③ 设备按指针差 DMA 取走命令
  \draw[hw bus] (sq4.north) -- (6.35,1.85) -- (10.45,1.85) -- (dev.north);
  \node[hw label, above] at (8.3,1.85) {③ DMA 取命令};
  % doorbell 通知设备（只携带新尾指针）
  \draw[hw return] (db.east) -- (dev.west);
  % ④ 设备写完成项进 CQ（device.south 竖直干线共享）
  \draw[hw bus] (dev.south) -- (10.45,-1.5) -- (6.35,-1.5) -- (cq4.south);
  \node[hw label, below] at (7.3,-1.5) {④ 写完成项};
  % ⑤ MSI 中断回本核
  \draw[hw return] (dev.south) -- (10.45,-2.6) -- (0.3,-2.6) -- (0.3,0.07);
  \node[hw label, below] at (5.0,-2.6) {⑤ MSI 中断回本核};
\end{pdfdiagram}
\diagramnote{NVMe 的提交—完成链路。驱动把 64 B 命令项写进本核 SQ 尾槽（①），屏障之后写 SQ 尾 doorbell（②）——doorbell 只携带新尾指针；设备按指针差 DMA 取走命令（③），执行后把 16 B 完成项写进 CQ（④），此时数据已按 PRP 地址 DMA 进入内存，最后以 MSI-X 中断通知本核（⑤）。SQ 与 CQ 各自独立、每核一对，提交路径上没有需要上锁的共享状态。}

第二处改动是多队列。规范允许最多 65535 对 I/O 队列，每对深度上限
65535；真实盘实现的对数远小于此（几十到几千），但已足够给每个 CPU
核配一对。每核一对的用意是免锁：各核写自己的 SQ、收自己的 CQ，尾指
针是核私有的，提交路径上没有任何共享状态需要上锁。这不是锁实现得好，
而是根本没有共享——“冒险、异常与性能”一章谈多核时说过，把共享拆开比把锁做快更值钱，
NVMe 的多队列是这句话在外设接口上的直接应用。配合 MSI-X 的多向量中
断，完成也能按队列分发回本核处理，提交与完成两条路径都不出核。

把每核一对的结构画出来，免锁与不出核是同一件事的两个侧面：

\begin{pdfdiagram}
  % 左：各核
  \node[hw compute, minimum width=1.1cm, minimum height=0.62cm] (k0) at (0.85,0.5) {核 0};
  \node[hw compute, minimum width=1.1cm, minimum height=0.62cm] (k1) at (0.85,-0.6) {核 1};
  \node[hw label] at (0.85,-1.55) {$\vdots$};
  \node[hw compute, minimum width=1.1cm, minimum height=0.62cm] (kn) at (0.85,-2.5) {核 $n$};
  % 中：每核一对 SQ/CQ（主机内存）
  \node[hw memory, minimum width=3.1cm] (p0) at (4.05,0.5) {SQ 0：命令环\\CQ 0：完成环};
  \node[hw memory, minimum width=3.1cm] (p1) at (4.05,-0.6) {SQ 1：命令环\\CQ 1：完成环};
  \node[hw label] at (4.05,-1.55) {$\vdots$};
  \node[hw memory, minimum width=3.1cm] (pn) at (4.05,-2.5) {SQ $n$：命令环\\CQ $n$：完成环};
  \node[hw compute, minimum width=2.6cm] (ctl) at (8.35,-1.0) {NVMe 控制器\\FTL＋DMA 引擎};
  % 提交：命令项 + doorbell
  \draw[hw bus] (k0.east) -- (p0.west);
  \draw[hw bus] (k1.east) -- (p1.west);
  \draw[hw bus] (kn.east) -- (pn.west);
  \node[hw label, anchor=south] at (1.95,1.0) {写命令项};
  \draw[BookMuted, line width=0.62pt] (6.15,0.5) -- (6.15,-2.5);
  \draw[hw bus] (p0.east) -- (6.15,0.5);
  \draw[hw bus] (p1.east) -- (6.15,-0.6);
  \draw[hw bus] (pn.east) -- (6.15,-2.5);
  \draw[hw bus] (6.15,-1.0) -- (ctl.west);
  \node[hw label, anchor=east] at (6.0,-1.15) {DMA 取命令};
  % doorbell：只带新尾指针
  \draw[hw return] (kn.south) -- (0.85,-3.35) -- (8.35,-3.35) -- (ctl.south);
  \node[hw label, anchor=north] at (4.6,-3.42) {doorbell：只带新尾指针（每队列一对寄存器）};
  % 完成：写 CQ + MSI-X 回本核
  \draw[hw return] (ctl.north) -- (8.35,1.35) -- (4.05,1.35) -- (p0.north);
  \node[hw label, anchor=south] at (6.2,1.42) {写 CQ 完成项，MSI-X 回本核};
  \node[hw label] at (4.6,-4.25) {尾指针核私有：提交路径上没有共享状态；规范上限 65535 对、每对深 65535};
\end{pdfdiagram}
\diagramnote{NVMe 的多队列结构：每核一对 SQ/CQ，尾指针核私有，doorbell 只携带新尾指针；设备经 DMA 取走命令、写 CQ 完成项，再以 MSI-X 按队列分发回本核——提交与完成两条路径都不出核。免锁不是锁实现得好，而是提交路径上根本没有共享状态。规范允许最多 65535 对队列、每对深 65535，真实盘实现几十到几千对，常用深度 32 至 1024。}

对照 AHCI，就能看出这两处改动解决了什么。AHCI 整台控制器只有一张命
令表，32 个槽位（NCQ 深度 32），一个完成中断源：所有核的提交都抢同
一张表的锁，完成中断也集中在一处。在 SATA 的 600 MB/s、约十万 IOPS
的上限下，这张表不是瓶颈；PCIe SSD 把带宽抬高一个数量级、把 IOPS
抬高两个数量级，提交路径上的共享状态先成为瓶颈。NVMe 的多队列不是
锦上添花，而是把 AHCI 那张唯一的表拆开之后的必然形态。

把两代接口的提交结构并排画出来，瓶颈的出处一目了然：

\begin{pdfdiagram}
  % 分隔线
  \draw[BookRule, line width=0.45pt, dashed] (5.35,0.9) -- (5.35,-3.0);
  % 左：AHCI
  \node[hw label] at (2.55,0.95) {AHCI：整控制器一张命令表（32 槽）};
  \node[hw compute, minimum width=0.9cm, minimum height=0.55cm] (c0) at (0.55,0.35) {核 0};
  \node[hw compute, minimum width=0.9cm, minimum height=0.55cm] (c1) at (0.55,-0.35) {核 1};
  \node[hw compute, minimum width=0.9cm, minimum height=0.55cm] (c2) at (0.55,-1.05) {核 2};
  \node[hw compute, minimum width=0.9cm, minimum height=0.55cm] (c3) at (0.55,-1.75) {核 3};
  \draw[BookMuted, line width=0.62pt] (1.35,0.35) -- (1.35,-1.75);
  \draw[hw bus] (c0.east) -- (1.35,0.35);
  \draw[hw bus] (c1.east) -- (1.35,-0.35);
  \draw[hw bus] (c2.east) -- (1.35,-1.05);
  \draw[hw bus] (c3.east) -- (1.35,-1.75);
  \node[hw control, minimum width=1.3cm] (lk) at (2.35,-0.7) {表锁};
  \node[hw memory, minimum width=2.0cm] (tb) at (4.1,-0.7) {命令表\\32 槽（NCQ）};
  \draw[hw bus] (1.35,-0.7) -- (lk.west);
  \draw[hw bus] (lk.east) -- (tb.west);
  \node[hw label] at (2.55,-2.6) {多核提交抢同一把表锁\\完成中断集中一处};
  % 右：NVMe
  \node[hw label] at (8.6,0.95) {NVMe：每核一对 SQ/CQ};
  \node[hw compute, minimum width=0.9cm, minimum height=0.55cm] (d0) at (6.25,0.35) {核 0};
  \node[hw compute, minimum width=0.9cm, minimum height=0.55cm] (d1) at (6.25,-0.35) {核 1};
  \node[hw compute, minimum width=0.9cm, minimum height=0.55cm] (d2) at (6.25,-1.05) {核 2};
  \node[hw compute, minimum width=0.9cm, minimum height=0.55cm] (d3) at (6.25,-1.75) {核 3};
  \node[hw memory, minimum width=1.9cm, minimum height=0.55cm] (q0) at (8.05,0.35) {SQ/CQ 0};
  \node[hw memory, minimum width=1.9cm, minimum height=0.55cm] (q1) at (8.05,-0.35) {SQ/CQ 1};
  \node[hw memory, minimum width=1.9cm, minimum height=0.55cm] (q2) at (8.05,-1.05) {SQ/CQ 2};
  \node[hw memory, minimum width=1.9cm, minimum height=0.55cm] (q3) at (8.05,-1.75) {SQ/CQ 3};
  \draw[hw bus] (d0.east) -- (q0.west);
  \draw[hw bus] (d1.east) -- (q1.west);
  \draw[hw bus] (d2.east) -- (q2.west);
  \draw[hw bus] (d3.east) -- (q3.west);
  \draw[BookMuted, line width=0.62pt] (9.2,0.35) -- (9.2,-1.75);
  \draw[hw bus] (q0.east) -- (9.2,0.35);
  \draw[hw bus] (q1.east) -- (9.2,-0.35);
  \draw[hw bus] (q2.east) -- (9.2,-1.05);
  \draw[hw bus] (q3.east) -- (9.2,-1.75);
  \node[hw compute, minimum width=1.7cm] (nc) at (10.25,-0.7) {NVMe\\控制器};
  \draw[hw bus] (9.2,-0.7) -- (nc.west);
  \node[hw label] at (8.6,-2.6) {核私有指针，提交免锁\\MSI-X 按核分发完成};
\end{pdfdiagram}
\diagramnote{AHCI 与 NVMe 的提交结构对照。AHCI 整台控制器只有一张 32 槽命令表（NCQ 深度 32）：所有核的提交抢同一把表锁，完成中断也集中在一处；SATA 的 600 MB/s、约十万 IOPS 下这张表不是瓶颈，PCIe SSD 把带宽抬高一个数量级、IOPS 抬高两个数量级之后，提交路径上的共享状态先成为瓶颈。NVMe 把表拆开：每核一对 SQ/CQ，尾指针核私有、提交免锁，完成由 MSI-X 按核分发——多队列是把那张唯一的表拆开之后的必然形态。}

\begin{longtable}{L{0.16\linewidth}L{0.30\linewidth}L{0.34\linewidth}}
\toprule
维度 & AHCI（SATA 时代） & NVMe（PCIe 时代） \\
\midrule
队列数 & 整控制器一张命令表 & 规范上限 65535 对 SQ/CQ \\\tablerowrule
每队深度 & 32（NCQ 槽位） & 上限 65535，常用 32 至 1024 \\\tablerowrule
提交路径 & 共享表，多核提交要抢锁 & 每核一对，核私有指针，免锁 \\\tablerowrule
完成通知 & 单一中断源 & MSI-X 多向量，可按核分发 \\\tablerowrule
典型上限 & 600 MB/s、约十万 IOPS & 数 GB/s、几十万到百万 IOPS \\
\bottomrule
\end{longtable}

深度上限 65535 也不意味着要用满。队列深度的用途是维持在途请求数，
“冒险、异常与性能”一章的 Little 定律给出该维持多少；用满上限买到的是排队而不是吞
吐——本节第四主题会把这笔账算到具体数字。

\topic{案例：一次 NVMe 读的完整 trace}
把一个进程执行 \code{read(fd, buf, 4096)}、且数据不在页缓存的全过
程按时序列出。设备取一块典型的 PCIe 3.0 x4 NVMe SSD，请求是 4 KB
随机读。

\tablechunkintro{1}{2}{1}{7}
\begin{longtable}{L{0.06\linewidth}L{0.44\linewidth}L{0.14\linewidth}L{0.16\linewidth}}
\toprule
步 & 动作 & 耗时量级 & 执行者 \\
\midrule
1 & 系统调用陷入内核，检查参数与文件描述符 & 约 1 $\mu$s & CPU \\\tablerowrule
2 & VFS 查页缓存：未命中，分配空闲页框 & 微秒级 & CPU \\\tablerowrule
3 & 文件系统把文件偏移翻译成设备 LBA & 微秒级 & CPU \\\tablerowrule
4 & 块层生成请求，调度器尝试与相邻请求合并 & 微秒级 & CPU \\\tablerowrule
5 & 驱动把请求编成 64 B 命令项，写入本核 SQ 尾槽 & 微秒级 & CPU \\\tablerowrule
6 & 屏障之后写 SQ 尾 doorbell（MMIO 写） & 百纳秒级 & CPU \\\tablerowrule
7 & 设备 DMA 取走命令项 & 微秒级 & 设备 \\
\bottomrule
\end{longtable}

\tablechunkintro{2}{2}{8}{13}
\begin{longtable}{L{0.06\linewidth}L{0.44\linewidth}L{0.14\linewidth}L{0.16\linewidth}}
\toprule
步 & 动作 & 耗时量级 & 执行者 \\
\midrule
8 & 查 FTL 映射，读 NAND 页，ECC 校验 & 约 80 $\mu$s & 介质 \\\tablerowrule
9 & 数据按 PRP 地址 DMA 写入第 2 步的页框 & 约 1 $\mu$s & 设备 \\\tablerowrule
10 & 设备写 16 B CQ 完成项，相位位翻转 & 微秒级 & 设备 \\\tablerowrule
11 & 设备发 MSI-X，中断控制器转成 CPU 中断 & 微秒级 & 设备 \\\tablerowrule
12 & 中断处理读 CQ，找回请求，标记完成，唤醒进程 & 微秒级 & CPU \\\tablerowrule
13 & 页缓存数据拷入用户缓冲，\code{read} 返回 4096 & 微秒级 & CPU \\
\bottomrule
\end{longtable}

总时延约 100 $\mu$s，其中第 8 步介质本身占去八成，软件与队列合计约
两成——这就是 NVMe 时代的比例：协议和驱动的开销已被压到介质时延的
零头，继续优化软件路径收益有限，优化介质的访问形态（对齐、合并、足
够的在途深度）才是主战场。这张表也是排障地图：慢在哪个量级，就查哪
一列——微秒级的步骤出问题多半是驱动或内核配置，毫秒级的问题才轮到
介质与垃圾回收。

第 6 步值得停一下：先写命令项、再写 doorbell，中间隔一道屏障，这正
是存储与设备事务相关章节“描述符先可见、再通知”的约定；顺序反了，设备可能读到半写
入的命令项。第 11 步也值得一提：MSI 不是一根专用电平线，而是设备向
中断控制器的特定地址做一次写事务。“电源、时钟与信号完整性”一章讨论过沿导线传播的信号要付
出走线与时序的代价；MSI 干脆把中断变成数据流里的一次写，省掉专用引
脚，还能在写数据里携带向量号，MSI-X 于是能按队列把中断分发到不同
核——中断从“一根线”演化成“一条消息”，和多队列是同一个设计方向。

同一件事在经典芯片与平台案例的 BIOS 路径上长什么样，对照一下就清楚了：

\begin{longtable}{L{0.14\linewidth}L{0.34\linewidth}L{0.34\linewidth}}
\toprule
维度 & int 13h（经典芯片与平台案例 BIOS 路径） & NVMe 路径 \\
\midrule
提交 & 参数装寄存器，执行 \code{int 0x13} & 命令项写内存队列，doorbell 通知 \\\tablerowrule
地址 & CHS 柱面/磁头/扇区 & LBA 加 PRP 物理页列表 \\\tablerowrule
等待 & BIOS 轮询状态寄存器，独占 CPU & 进程睡眠，完成由 MSI 唤醒 \\\tablerowrule
并发 & 一次一个在途 & 几百在途，多队列并行 \\\tablerowrule
搬运 & 软盘走 8237 DMA，硬盘 PIO 或总线主控 & 设备直接 DMA 到页框 \\
\bottomrule
\end{longtable}

骨架是同一个：准备命令、通知设备、DMA 搬数、确认完成。int 13h 是这
套骨架的同步简化版——提交即等待，全程没有队列，CPU 停在 BIOS 例程
里轮询；NVMe 是它的并发完整版——提交与完成解耦，等待换成中断，在
途请求从一个变成几百个。从 8237 的一个 DMA 通道到每核一对队列，外
设接口四十年的演进主线，就是把“一次一单”改成“持续事务流”，而描
述符与 DMA 这两个零件自始至终没有换过。

\topic{尾延迟：P99 从哪里来}
前面反复出现的“典型时延 100 $\mu$s”只是中位数附近的值。“冒险、异常与性能”一章已
经立下规矩：时延是一个分布，平均值会撒谎。外设是把这条规矩演绎得最
充分的地方——慢请求不是随机噪声，而是有明确物理来源的事件，每一类
都对应本章前几节讲过的机制：

\begin{longtable}{L{0.16\linewidth}L{0.36\linewidth}L{0.14\linewidth}L{0.20\linewidth}}
\toprule
长尾来源 & 物理机制 & 典型幅度 & 触发条件 \\
\midrule
GC 擦除 & 写请求排在一次块擦除与有效页搬运后面 & 毫秒级 & 空闲块耗尽、盘接近写满 \\\tablerowrule
队列拥塞 & 深队列高负载下排队，“冒险、异常与性能”一章的 $1/(1-\rho)$ & 百微秒级 & 深度过大或负载接近饱和 \\\tablerowrule
SMR 重写 & 区域写满触发整段重写（第三节） & 十毫秒级至秒级 & 随机写命中已写区域 \\\tablerowrule
读重试 & 阈值电压漂移，ECC 纠不动，换挡重读 & 数十微秒至毫秒 & 老盘、高温、久置数据 \\\tablerowrule
映射回读 & FTL 映射在 DRAM 中未命中，回读 NAND 映射页 & 百微秒级 & 冷地址、大跨度跳跃 \\
\bottomrule
\end{longtable}

这些事件不常发生，所以中位数很好看；但一旦发生就是十倍、百倍的慢。
消费级 NVMe SSD 上 4 KB 随机读的典型分布是：

\begin{longtable}{L{0.14\linewidth}L{0.16\linewidth}L{0.20\linewidth}L{0.32\linewidth}}
\toprule
百分位 & 时延 & 相对中位数 & 直观含义 \\
\midrule
P50 & 100 $\mu$s & 1 倍 & 典型请求 \\\tablerowrule
P99 & 1 ms & 10 倍 & 每 100 个里约 1 个这么慢 \\\tablerowrule
P999 & 10 ms & 100 倍 & 每 1000 个里约 1 个这么慢 \\
\bottomrule
\end{longtable}

把这张分布画成直方图，长尾的形状比数字更直接：

\begin{pdfdiagram}
  % 坐标轴
  \draw[BookMuted, line width=0.6pt, ->] (0.5,0) -- (11.0,0);
  \draw[BookMuted, line width=0.6pt, ->] (0.5,0) -- (0.5,4.0);
  \node[hw label, anchor=south] at (0.5,4.05) {请求数};
  \node[hw label] at (6.9,-0.3) {时延（横轴按对数间距）};
  % 直方图：主峰在 100 µs 附近，长尾拖向 1 ms、10 ms
  \filldraw[fill=BookTeal!25, draw=BookTeal, line width=0.5pt] (0.8,0) rectangle (1.3,0.7);
  \filldraw[fill=BookTeal!25, draw=BookTeal, line width=0.5pt] (1.35,0) rectangle (1.85,3.2);
  \filldraw[fill=BookTeal!25, draw=BookTeal, line width=0.5pt] (1.95,0) rectangle (2.45,1.5);
  \filldraw[fill=BookTeal!25, draw=BookTeal, line width=0.5pt] (2.55,0) rectangle (3.05,0.8);
  \filldraw[fill=BookTeal!25, draw=BookTeal, line width=0.5pt] (3.25,0) rectangle (3.75,0.5);
  \filldraw[fill=BookTeal!25, draw=BookTeal, line width=0.5pt] (4.05,0) rectangle (4.55,0.35);
  \filldraw[fill=BookTeal!25, draw=BookTeal, line width=0.5pt] (5.05,0) rectangle (5.55,0.25);
  \filldraw[fill=BookTeal!25, draw=BookTeal, line width=0.5pt] (6.25,0) rectangle (6.75,0.18);
  \filldraw[fill=BookTeal!25, draw=BookTeal, line width=0.5pt] (7.65,0) rectangle (8.15,0.13);
  \filldraw[fill=BookTeal!25, draw=BookTeal, line width=0.5pt] (9.15,0) rectangle (9.65,0.1);
  \filldraw[fill=BookTeal!25, draw=BookTeal, line width=0.5pt] (10.15,0) rectangle (10.65,0.07);
  % 百分位参考线与平均值
  \draw[BookMuted, line width=0.5pt, dashed] (1.6,0) -- (1.6,3.75);
  \node[hw label, anchor=west] at (1.72,3.9) {P50 = 100 $\mu$s};
  \draw[BookRed, line width=0.5pt, dashed] (2.5,0) -- (2.5,3.35);
  \node[hw label, anchor=west, text=BookRed] at (2.62,3.5) {平均值 $\approx$199 $\mu$s};
  \draw[BookRed, line width=0.6pt, ->] (1.72,2.9) -- (2.4,2.9);
  \node[hw label, anchor=west] at (2.62,2.9) {被长尾拉偏};
  \draw[BookMuted, line width=0.5pt, dashed] (9.4,0) -- (9.4,1.5);
  \node[hw label, anchor=south] at (9.4,1.55) {P99 = 1 ms};
  \draw[BookMuted, line width=0.5pt, dashed] (10.4,0) -- (10.4,1.1);
  \node[hw label, anchor=south] at (10.4,1.15) {P999 = 10 ms};
  % 横轴刻度与长尾来源
  \node[hw label, anchor=north] at (1.6,-0.12) {100 $\mu$s};
  \node[hw label, anchor=north] at (9.4,-0.12) {1 ms};
  \node[hw label, anchor=north] at (10.4,-0.12) {10 ms};
  \node[hw label] at (6.2,0.85) {长尾：GC 擦除、队列拥塞、SMR 重写、读重试};
\end{pdfdiagram}
\diagramnote{消费级 NVMe SSD 4 KB 随机读的时延分布（横轴按对数间距）。主峰在 P50 = 100 $\mu$s；GC 擦除、队列拥塞、SMR 重写、读重试这类不常发生、一发生就是十倍百倍慢的事件，把分布拖出长尾：P99 = 1 ms，P999 = 10 ms。平均值被长尾向右拉偏，离中位数越来越远——所以外设时延必须带百分位一起报。}

平均时延把这些全藏了起来。算一遍“冒险、异常与性能”一章的同款例子：100 个请求里 99
个花 100 $\mu$s、1 个花 10 ms，平均值是
$(99\times100+10000)/100=199\ \mu$s——平均看上去只比中位数差一倍，
可那一个请求慢了 100 倍。若服务目标是“99\% 的请求快于
200 $\mu$s”，一份只报平均值 199 $\mu$s 的报告，会把一次彻底的违约
藏进“接近达标”的表象里。用户记住的从来不是平均那次，而是最慢那
次：页面偶尔卡住三秒，原因往往就是排在 P999 上的那个请求撞上了 GC
擦除或 SMR 重写。

工程含义有三条。其一，报外设时延必须带百分位：平均值至少配 P99，交
互与实时负载再看 P999，只报平均值的测试报告不可信。其二，容量规划
按 P99 反推负载上限：让 P99 保持在目标内的负载，通常只有平均时延开
始明显变坏时的一半左右——这与“冒险、异常与性能”一章“稳态利用率压在五到七成”是同
一条纪律在外设上的形式。其三，长尾不能靠平均优化消掉，只能按来源逐
个治理：GC 尖峰靠预留空闲块和 TRIM，队列拥塞靠限深或限流，SMR 重写
靠改变写入形态——这正是本节后面两个主题的内容。

\topic{队列深度与吞吐时延的权衡}
“冒险、异常与性能”一章的 Little 定律 $L=\lambda W$ 用在外设上的形式是：在途请求数
等于 IOPS 乘单请求时延。介质时延 $W$ 由物理决定（4 KB 随机读约
100 $\mu$s），吞吐目标 $\lambda$ 由应用决定，队列深度是主机侧唯一
的旋钮。两个方向都有代价。

深度太小，喂不饱设备。QD=1 时 $\lambda=1/W=10^{4}$ IOPS，按 4 KB
折算只有 40 MB/s——一块标称 3.5 GB/s 的盘只用了约 1\%。深度加大，
吞吐近似线性上升，直到撞上设备自身的并行度上限（NAND 通道数与控制
器能力）。设某盘上限为 50 万 IOPS，拐点的位置由
$L=\lambda W=5\times10^{5}\times10^{-4}=50$ 唯一决定。越过拐点再加
深度，吞吐不变，多出来的在途全部变成排队：QD=256 时平均时延约
$256/(5\times10^{5})=512\ \mu$s，是介质时延的五倍多。

\begin{longtable}{L{0.12\linewidth}L{0.18\linewidth}L{0.18\linewidth}L{0.34\linewidth}}
\toprule
队列深度 & 可达 IOPS & 平均时延 & 位置 \\
\midrule
1 & 1 万 & 100 $\mu$s & 深度限制，带宽饿着 \\\tablerowrule
4 & 4 万 & 100 $\mu$s & 深度限制 \\\tablerowrule
16 & 16 万 & 100 $\mu$s & 深度限制 \\\tablerowrule
32 & 32 万 & 100 $\mu$s & 接近拐点 \\\tablerowrule
64 & 50 万（封顶） & 约 128 $\mu$s & 刚过拐点，开始排队 \\\tablerowrule
128 & 50 万 & 约 256 $\mu$s & 排队区，吞吐不再涨 \\\tablerowrule
256 & 50 万 & 约 512 $\mu$s & 排队区，纯买时延 \\
\bottomrule
\end{longtable}

这条曲线的形状与“冒险、异常与性能”一章的 roofline 是同一族：拐点左侧被深度（在途
数）限制，右侧被设备能力封顶，拐点本身由 $\lambda W$ 决定。工程读
法也只有一句：深度按 $\lambda W$ 取，再留一档余量；过了拐点一分钱
吞吐也买不到，买到的全是时延。

主机侧的经验值：每队列深度取 32 至 128 是常见区间；\code{fio} 压测
的惯例 \code{iodepth=32} 正落在典型拐点附近——这是压测惯例，不是
默认值：\code{fio} 默认走同步 I/O，\code{iodepth=1}（本章末尾的案
例一就是踩在这个默认值上）。时延敏感的负载反而要浅：日志刷盘、同步写这类请求应走独立队列
或浅队列，避免排在批量流量后面——用“冒险、异常与性能”一章的语言说，浅队列买的是时
延的可预期性，不是平均速度。多队列改变的是这笔账的分母：$n$ 个核各
开一对队列，每对只承担 $\lambda/n$，同样的总吞吐用更浅的每队深度就
够，排队时延同步下降——NVMe 的多队列与深度取舍，是同一条定律的两
个应用。

\topic{调试案例两则}
两则案例的共同点是：症状都出现在吞吐或时延上，根因都在队列机制与介
质特性的交界处，定位方法都是“先看队列计数，再看介质状态”。

案例一：队列深度 1 的吞吐瓶颈。症状：新服务器的 NVMe SSD 规格
3.5 GB/s，实测顺序读只有约 40 MB/s，4 KB 随机读约 1 万 IOPS，离规格
差两个数量级。定位：\code{iostat -x} 显示设备利用率 100\%，但队列
长度 \code{aqu-sz} 恒为 1.00；压测用的是 \code{fio} 默认的同步
I/O，\code{iodepth=1}。根因：同步 I/O 一次只有一个请求在途，
Little 定律直接封顶——$4\ \text{KB}/100\ \mu\text{s}=40$ MB/s。设
备没有病，是提交方式只给它一单做一单。修复：改用 \code{libaio} 或
\code{io\_uring}，\code{iodepth} 提到 32。同机复测：4 KB 随机读从 1
万 IOPS 升到约 32 万 IOPS，顺序读回到 3 GB/s 量级——吞吐涨了约 32
倍，硬件一分钱没动。教训：报外设性能必须先报队列深度，没有深度的
IOPS 数字没有意义；看到利用率 100\% 先看队列长度——利用率满而队列
空，瓶颈在提交侧，不在设备侧。

案例二：GC 风暴导致的周期性延迟尖峰。症状：持续随机写的服务，P50
写时延稳定在 200 $\mu$s，但每隔几秒出现一批 10 至 50 ms 的慢写，
P99 从 1 ms 恶化到 30 ms；写入速率越快，尖峰间隔越短。定位：时延直
方图呈双峰；盘已写到接近满容量；SSD 健康信息里主机写入量远小于
NAND 实际写入量，写放大超过 3。根因：盘写满后没有现成空闲块，每次
写都要先做一次 GC——读出整块、搬走有效页、擦除——写请求排在一次
毫秒级的擦除后面；空闲块耗尽一批就尖峰一次，所以呈周期性，且写入越
快耗尽越快。对照实验：对盘做安全擦除（或全量 TRIM）后重测，尖峰消
失、P99 回落，确认根因在空闲块供给，不在队列深度或接口。修复：文件
系统开启 \code{discard} 或定期 \code{fstrim}，让删除的文件及时变回
设备可用的空闲块；预留空间提到 20\% 以上；时延敏感的服务不把盘写
满，容量水位线设在 70\%。用闲置容量换空闲块的持续供给，本质上是给
GC 留余量——与“冒险、异常与性能”一章“服务台保持余量”是同一条纪律。

两则案例放在一起看，方法上的共性比结论本身更值得带走。两案的设备都
是完好的，换硬件解决不了任何一案；两案的突破口都是计数器而不是猜
测——队列长度、利用率、写放大、时延直方图，全是现成的可观测数据；
两案的修复也都不动设备，一个改提交方式，一个改介质水位。外设问题的
默认值应当是“机制问题”而不是“硬件坏了”：先把 Little 定律和介质的
物理与协议约束代入算一遍，算不通再怀疑硬件。

\topic{外设问题的通用读法}
全章收束为一张检查表。面对任何外设的性能或异常问题，按三步走：先问
介质特性——这类设备的物理限制是什么；再问队列机制——它用什么结构
弥补介质的短板；最后测量验证——不看规格书看计数器，不看平均看百分
位。三类设备放进同一张表：

\begin{longtable}{L{0.10\linewidth}L{0.26\linewidth}L{0.26\linewidth}L{0.26\linewidth}}
\toprule
设备 & 介质先问 & 队列再问 & 测量验证 \\
\midrule
SSD & 空闲块剩多少，写放大多少 & 提交深度多少，谁在排队 & \code{fio} 按百分位报时延，\code{iostat} 看 \code{aqu-sz} \\\tablerowrule
机械盘 & 随机还是顺序，寻道占几成 & 调度器合并了吗，NCQ 用上了吗 & \code{iostat} 看 await 与服务时间之差 \\\tablerowrule
网卡 & 先到 pps 顶还是带宽顶，帧多大 & 聚合参数多少，多队列吃满核了吗 & \code{ethtool -S} 看丢包计数，按核看中断分布 \\
\bottomrule
\end{longtable}

三类设备共享同一张检查表不是巧合：存储与设备事务相关章节的描述符环与 DMA、“冒险、异常与性能”一章的
Little 定律与百分位纪律，在每一块外设上都是同一套语言。读外设问题
的功力，说到底就是先认出介质那一栏，再走完后两栏。

把三步走完整演练一次。某数据库机器报告“磁盘慢”：第一步问介质——这
是一块 SATA SSD 还是 NVMe SSD、写到了几成满、随机写占多大比例，因
为满盘加随机写正是 GC 尖峰的经典触发条件；第二步问队列——提交深度
多少、走的是单队列还是多队列、同步写有没有和批量读混在同一对队列
里，因为平均时延 200 $\mu$s 但 P99 到 20 ms 的形状，指向的是“排在
慢事件后面”而不是“介质本身慢”；第三步测量验证——\code{fio} 单独复
现负载并按百分位报时延，\code{iostat -x} 对照 \code{aqu-sz} 与利用
率，再查 SSD 健康信息里的写放大与剩余备用块。三步走完，“磁盘慢”这
句模糊报告就被翻译成一句可执行的话：盘已到 92\% 容量、写放大 4.5，
先把水位降到 70\% 再谈别的。同一个流程搬到网卡上，只是第一栏换成
pps 与帧长、第三栏换成 \code{ethtool} 计数——方法不动，参数换设备。




























\section{三、以太网 PHY：从 MAC 帧到介质信号}

网卡驱动建立描述符并启动 DMA 后，MAC 得到的是一串帧字节；线缆或光纤传输的却是经过编码、加扰、纠错和模拟发送的符号。二者之间由协调子层、介质无关接口、PCS、PMA 和 PMD/MDI 等层次连接。不同速率和介质采用的编码与训练细节不同，但从 MAC 到 PHY 的责任边界保持一致：MAC 管理帧和链路层地址，PHY 管理符号、时钟恢复、信道适配与介质电气/光学接口。\cite{ieee8023-baselines}

\topic{发送路径逐层改变表示形式}
发送时，MAC 先形成前导码、帧头、有效载荷和 FCS，并通过 MII 家族接口把数据交给 PHY。PCS 完成编码、加扰、lane 分配与对齐标记；需要前向纠错时，FEC 编码器加入冗余；PMA 负责串并转换、lane 复用和时钟关系；PMD 最终驱动铜缆、背板或光模块。每层都可能扩大数据量或改变时钟域，因此“MAC 每秒交出多少有效载荷”不能直接等同于“介质上需要多少符号带宽”。

接收方向按相反顺序推进。模拟前端先均衡并判决采样，时钟数据恢复电路从信号中恢复采样相位，PMA 形成稳定符号流，FEC 尝试纠错，PCS 解码并恢复 lane 顺序，MAC 最后校验 FCS 并决定接收、统计错误或丢弃。一个 FCS 错误只说明帧到达 MAC 时内容不一致；要继续定位，还需查看纠错计数、符号错误、lane 对齐、信噪比和收发光功率等 PHY 证据。

\begin{pdfdiagram}[x=1cm,y=1cm]
  % 发送路径和接收路径分开占两排，避免返回箭头穿过节点。
  \node[hw label, font=\small\bfseries] at (-0.15,3.85) {发送};
  \node[hw state, minimum width=2.05cm, align=center] (txdma) at (1.05,3.2) {主机内存\\TX 描述符与帧};
  \node[hw control, minimum width=2.05cm, align=center] (txmac) at (3.8,3.2) {MAC\\地址、长度、FCS};
  \node[hw compute, minimum width=2.05cm, align=center] (txpcs) at (6.55,3.2) {PCS / FEC\\编码与加扰};
  \node[hw compute, minimum width=2.05cm, align=center] (txpma) at (9.3,3.2) {PMA / PMD\\串行化与驱动};
  \draw[hw bus] (txdma) -- (txmac);
  \draw[hw bus] (txmac) -- (txpcs);
  \draw[hw bus] (txpcs) -- (txpma);

  \node[hw label, font=\small\bfseries] at (-0.15,0.95) {接收};
  \node[hw state, minimum width=2.05cm, align=center] (rxdma) at (1.05,0.3) {RX completion\\DMA 到主机页};
  \node[hw control, minimum width=2.05cm, align=center] (rxmac) at (3.8,0.3) {MAC\\FCS 与帧过滤};
  \node[hw compute, minimum width=2.05cm, align=center] (rxpcs) at (6.55,0.3) {PCS / FEC\\纠错与解码};
  \node[hw compute, minimum width=2.05cm, align=center] (rxpma) at (9.3,0.3) {PMA / PMD\\采样与时钟恢复};
  \draw[hw return] (rxpma) -- (rxpcs);
  \draw[hw return] (rxpcs) -- (rxmac);
  \draw[hw return] (rxmac) -- (rxdma);

  \node[hw memory, minimum width=2.15cm, minimum height=1.2cm, align=center] (medium) at (11.85,1.75) {介质\\铜缆 / 光纤};
  \draw[hw bus] (txpma.east) -- (10.7,3.2) -- (10.7,2.25) -- (medium.north west);
  \draw[hw return] (medium.south west) -- (10.7,1.25) -- (10.7,0.3) -- (rxpma.east);
\end{pdfdiagram}
\diagramnote{一帧以不同表示穿过以太网接口。主机内存保存字节和描述符，MAC 处理帧语义，PCS/FEC 处理码块与纠错，PMA/PMD 处理 lane、采样和介质信号；接收完成只有在 PHY 恢复数据、MAC 接受帧并把内容 DMA 到主机页后才成立。每层都留下不同的错误证据。}

\topic{1 Gb/s 数值案例：64 B 最小帧并不只占 512 ns}

以不带 VLAN 标签的经典以太网帧为例，从目的 MAC 地址到 FCS 的最小帧长为 64 B，其中有效载荷不足时要填充到 46 B。介质上还要发送 7 B 前导码和 1 B SFD，并在下一帧前留出相当于 12 B 的帧间隔。因此，一次最小帧占用的线速时间为
\[
(8+64+12)\times8\ \mathrm{bit}\div10^9\ \mathrm{bit/s}
=672\ \mathrm{ns}.
\]
一个 1518 B 普通最大帧对应
\[
(8+1518+12)\times8\ \mathrm{bit}\div10^9\ \mathrm{bit/s}
=12.304\ \mu\mathrm{s}.
\]
这两个数字解释了为什么“每秒帧数”和“每秒有效载荷字节数”不是同一个指标：短帧中前导码、帧头、FCS 与帧间隔所占比例更大；驱动、DMA 和中断还会增加不在线路字节数中的主机开销。\cite{ieee8023-baselines}

\begin{longtable}{L{0.18\linewidth}L{0.23\linewidth}L{0.29\linewidth}L{0.20\linewidth}}
\toprule
边界 & 表示形式 & 本层确认什么 & 本层不能确认什么 \\
\midrule
TX 描述符 & 地址、长度、offload 与所有权位 & 驱动已经发布一块可 DMA 缓冲区 & 帧已经离开 PHY \\\tablerowrule
MAC & 前导、地址、类型/长度、payload、FCS & 帧格式成立并完成 FCS 生成或检查 & 模拟眼图和 FEC 裕量 \\\tablerowrule
PCS / FEC & 码块、加扰数据、对齐标记与冗余 & 码块同步，错误在能力范围内已纠正 & 远端应用已经消费数据 \\\tablerowrule
PMA / PMD & 串行符号、采样相位与电气/光学幅度 & lane 锁定并满足当前接收判决 & MAC 帧地址与上层校验正确 \\\tablerowrule
RX completion & 缓冲区、长度、状态与校验结果 & NIC 已把被接受的帧写入主机内存 & 应用线程已经运行并读取它 \\
\bottomrule
\end{longtable}

\topic{自动协商先确定共同能力}
链路两端上电后先交换能力，选择双方共同支持的速率、双工方式以及适用的流控或 FEC 参数。某些高速铜缆和背板链路还要执行链路训练：发送端提供规定序列，接收端测量信道响应并请求调整发射均衡，直到误码和裕量满足进入数据态的条件。协商完成只表示参数达成一致；训练、块锁定、lane 对齐和错误率稳定后，链路才真正可承载 MAC 流量。

排障时应把状态分成三层：没有协商结果，优先检查介质、对端能力和管理配置；已经协商但训练失败，检查损耗、端接、均衡和 lane；链路已进入数据态但错误持续增长，再检查 FEC 纠错余量、温度、连接器与串扰。直接重启驱动可能暂时重新训练，却不能证明物理问题消失。

\topic{FEC 用确定延迟换取误码容限}
假设接收端在一个 FEC 码块内检测到少量错误，译码器可利用冗余恢复原始数据并增加“已纠正码字”计数；错误超过纠错能力时，只能报告不可纠正码字，上层帧校验通常也会失败。FEC 改善的是给定信道条件下的有效误码率，不会增加模拟眼图本身的裕量，并且会引入编码冗余和固定流水延迟。

因此健康链路也可能持续出现少量已纠正错误。判断是否退化要观察单位时间纠错率、不可纠正错误、链路重训练次数和环境变化，而不是看到非零计数就立即更换硬件。若纠错率随温度或线缆弯折显著上升，它就是物理裕量收窄的早期证据。

\topic{从收帧异常回溯到物理层}
当应用报告吞吐下降时，先确认 MAC 是否出现丢帧、暂停帧或环耗尽；若 MAC 错误不多，再检查驱动队列和 CPU 调度。若 FCS、符号或 FEC 计数同步上升，则沿 PCS、PMA、PMD 回溯。这个顺序能避免把物理链路误判成软件队列问题，也能避免在 RX 环已经耗尽时反复调整模拟均衡。完整证据链应同时保存协商结果、训练状态、PHY 计数器、MAC 计数器、驱动队列和应用时间线。

\topic{案例：吞吐下降但 link 仍为 up}

假设应用吞吐从 940 Mb/s 降到 300 Mb/s，而管理界面仍显示 1 Gb/s、全双工。第一种可能是主机侧 RX 环耗尽：PHY 与 MAC 计数正常，但 \code{rx\_no\_buffer}、丢包和 CPU 软中断积压上升，此时应增加队列处理能力或修正中断亲和性。第二种可能是介质裕量下降：FCS、符号错误或已纠正 FEC 码字随时间增加，并伴随重训练，此时应检查线缆、连接器、光功率、温度和均衡。第三种可能是流控：收到大量 pause 帧时，发送端会按协议停顿，链路不掉线却无法持续发满。

定位的关键不是先重启驱动，而是把同一时间窗内的应用速率、RX/TX 环水位、MAC FCS 计数、PHY 符号/FEC 计数、pause 帧和重训练次数对齐。队列水位先恶化而 PHY 计数稳定，根因更接近主机；PHY 错误先增加而队列随后空闲，根因更接近介质。只有最早变化的证据才能指向最早失效边界。
\section{四、USB、HID 与音频接口的共同边界}

USB 是主机调度的分层总线。设备接入后先完成复位和枚举，主机读取描述符、分配地址、选择配置，再由驱动按 endpoint 类型提交传输。设备不能像 PCIe 主设备那样任意 DMA 系统内存，而是由主控制器按调度表发起总线事务，再由主控制器 DMA 主机缓冲区。

\begin{longtable}{L{0.18\linewidth}L{0.32\linewidth}L{0.4\linewidth}}
\toprule
传输类型 & 典型设备 & 优先保证 \\
\midrule
control & 枚举与配置 & 正确、可重试、双向短消息 \\\tablerowrule
bulk & U 盘、打印机 & 正确性和剩余带宽利用率 \\\tablerowrule
interrupt & 键盘、鼠标、触控 & 有界轮询间隔，不等同于物理中断线 \\\tablerowrule
isochronous & 音频、摄像头 & 周期带宽和时延，允许按协议处理丢失 \\
\bottomrule
\end{longtable}

HID 报告把按键或指针状态编码成固定格式；音频接口还要处理生产者与消费者时钟略有偏差的问题，靠反馈端点、缓冲水位或采样率转换避免缓冲区逐渐溢出/耗尽。传感器接口也遵循同一思路：先确定采样时钟和数据格式，再定义缓冲、时间戳、丢样与错误报告。

\section{五、从物理事件到进程可见事件}

一次键盘按键会经历开关抖动、设备端扫描与消抖、USB 报告、主控制器 DMA、完成中断、内核输入子系统和用户态读取。每层都可能合并或延迟事件，因此“按键发生时刻”至少有物理接触、设备采样、主机收包和进程读取四个时间戳。

排错时按层观察：设备是否枚举，endpoint 是否配置，主控制器队列是否前进，DMA 缓冲是否更新，完成中断是否到达，驱动是否把报告交给上层。只在应用程序里反复重试，无法区分线缆、电源、协议、控制器和驱动故障。

\section{六、USB 拓扑、枚举与描述符}

USB 的树形拓扑由主控制器、root hub、外部 hub 和设备组成。连接检测后，端口先完成供电与复位，设备在默认地址响应；主机读取设备描述符的基本字段，分配唯一地址，再读取完整的配置、接口和 endpoint 描述符。一个物理设备可包含多个 interface，例如摄像头同时提供视频、麦克风和控制接口，由不同类驱动绑定。

描述符是一组带长度与类型的记录。解析器必须验证长度、层级归属和 endpoint 数量，不能因为设备声称总长度很大就越界读取。选择 configuration 或 alternate setting 会改变设备启用的 endpoint 与带宽需求；高速周期传输若没有足够调度预算，配置阶段就应失败，而不是运行后随机丢包。

hub 不只是被动分线器。它报告端口连接、电源、过流和复位状态，并转发不同速度设备的事务。共享上行链路意味着多个 endpoint 竞争同一周期预算；逻辑上独立的设备仍可能共享 hub、电源和主控制器故障域。

\section{七、xHCI 队列与一次 USB 传输}

xHCI 主控制器使用内存中的 ring 描述工作。驱动把 TRB 写入 transfer ring，执行必要的 DMA 同步，再写 doorbell。控制器取走 TRB，通过相应 endpoint 发起总线事务，把数据 DMA 到缓冲区或从缓冲区读出，随后在 event ring 写入完成事件并触发中断。驱动根据 slot、endpoint 和 TRB 指针归属完成，再推进 ring。

ring 的 cycle bit 用于区分环回后的新旧条目。生产者只有在完整写好 TRB 后才交出所有权；消费者不能跨过 cycle 不匹配的条目。取消传输或复位 endpoint 时，还要处理已经预取的 TRB 和迟到事件。仅清空软件链表而不停止控制器，会让旧 DMA 写入已经复用的缓冲区。

\begin{pdfdiagram}[
  node distance=7mm and 8mm,
  usbbox/.style={draw, rounded corners, align=center, minimum width=24mm, minimum height=10mm, fill=blue!5},
  >=Stealth]
\node[usbbox] (driver) {驱动填写\\TRB 与缓冲};
\node[usbbox, right=of driver] (doorbell) {doorbell\\交出所有权};
\node[usbbox, right=of doorbell] (host) {xHCI 调度\\并执行 DMA};
\node[usbbox, right=of host] (usb) {USB 事务\\endpoint 响应};
\node[usbbox, below=of host] (event) {event ring\\完成码};
\node[usbbox, left=of event] (consume) {驱动回收\\缓冲与 TRB};
\draw[->] (driver) -- (doorbell);
\draw[->] (doorbell) -- (host);
\draw[->] (host) -- (usb);
\draw[->] (usb) |- (event);
\draw[->] (event) -- (consume);
\draw[->] (consume) -- node[left, font=\small]{新请求} (driver);
\end{pdfdiagram}
\diagramnote{USB 请求从驱动到设备再返回。doorbell 只交出描述符所有权，真正完成要等设备事务、DMA 与 event ring 状态全部闭合。}

\section{八、HID：报告、消抖与输入语义}

HID report descriptor 描述字段的用途、位宽、逻辑范围、单位和 report ID。键盘可以报告当前按下的键集合，鼠标可以报告相对位移，触控设备则包含多个 contact 的坐标与状态。驱动必须按 bit 字段解析，而不能假定每个值都按整字节对齐。

机械按键在接触瞬间会抖动，设备固件通过时间窗或状态机滤除快速翻转。之后主机按 interrupt endpoint 的轮询间隔获取报告。“interrupt”表示主机为它保留周期性轮询机会，并非设备直接拉起 CPU 中断线；主控制器完成一批传输后才通过平台中断通知 CPU。

输入系统还要处理键位按下、保持、重复和释放。若丢失包含释放状态的报告，软件可能认为按键一直按下；协议、设备或上层需要用完整状态报告、超时或重新同步恢复。安全相关输入还应区分设备身份与注入来源。

\section{九、音频流：时钟域、反馈与缓冲水位}

USB 音频把连续采样分装进周期传输。即使主机和设备都配置为 48 kHz，它们的晶振仍有微小频差：生产稍快会使缓冲逐渐上升，消费稍快会使缓冲逐渐下降。仅靠一个大缓冲只能延后溢出或欠载，不能消除长期速率差。

异步音频 endpoint 可用 feedback 告知主机实际消费速率，主机据此调整每个服务间隔发送的样本数。另一种系统使用采样率转换吸收频差。控制环要平滑噪声和抖动，调节过快会把水位波动变成可闻调制，调节过慢又可能越过缓冲边界。

音频时延由采样周期、USB 调度、设备缓冲、主控制器与软件队列共同组成。减小所有缓冲可以降低平均时延，却提高对调度抖动的敏感度。评价应同时报告端到端时延、抖动、欠载/溢出次数和时钟恢复状态。




























\section{十、传感器把连续世界变成带时间的数据}

传感器路径不是“读一个寄存器”这么简单。温度、压力、加速度或光强先被敏感元件转换为电压、电流、电容或电阻变化，模拟前端完成偏置、放大与滤波，模数转换器（Analog-to-Digital Converter, ADC）在采样时刻量化输入，数字接口再把样本、状态和时间信息交给控制器。驱动最终看到的整数同时受量程、分辨率、参考电压、校准系数和采样时刻约束。

若一个 12-bit ADC 的参考范围为 0--3.3 V，理想最低有效位约为 $3.3/4096\approx0.806$ mV。读数 2048 只说明输入接近中点；在换算物理量前还要应用零点偏移、增益误差和传感器传递函数。分辨率也不等于精度：噪声、参考源误差、非线性和板级耦合都可能使多个相邻码无法区分真实输入。

采样率必须覆盖信号带宽并留出模拟抗混叠滤波的过渡带。高于 Nyquist 频率的输入会折叠到低频，软件事后无法从样本中判断它原来位于哪个频段。提高数字采样率可以放宽模拟滤波器设计，却会增加总线、FIFO、DMA 和处理器预算；因此采样率是从模拟边界一直影响到内存带宽的系统参数。

\section{十一、I2C、SPI 与采样 FIFO 的边界}

低速配置常用 I2C：控制器按地址选择目标，通过寄存器号读写量程、采样率和中断阈值；目标可用 NACK 表示未接受，也可通过 clock stretching 延长某一位时间。较高速的样本流常用 SPI：片选界定目标和帧，时钟极性/相位决定采样边沿，全双工移位意味着“读寄存器”往往也要同时发送命令或空字节。两者都只保证接口上的 bit/byte 转移，不自动说明样本对应哪个物理时刻。

周期采样设备通常在内部定时器到期时锁存 ADC 结果并写入 FIFO。FIFO 水位达到阈值后产生中断，或由 DMA 批量搬入内存环形缓冲。控制器读取 FIFO 的动作晚于真实采样时刻；若应用需要跨设备对齐，时间戳应尽量靠近采样触发点生成，并说明属于传感器时钟、控制器时钟还是系统时钟。仅用中断到达 CPU 的时间作为样本时间，会把总线等待、中断合并和调度抖动一起计入。

\begin{longtable}{L{0.18\linewidth}L{0.30\linewidth}L{0.32\linewidth}L{0.12\linewidth}}
\toprule
边界 & 保存的状态 & 完成条件 & 常见故障 \\
\midrule
模拟采样 & 量程、参考、滤波与采样沿 & sample/hold 已锁定输入 & 饱和、混叠、参考噪声 \\\tablerowrule
ADC 转换 & 通道号、转换序号、校准参数 & 数字码与状态位有效 & 未等待转换完成 \\\tablerowrule
串行接口 & 设备地址/片选、寄存器、帧长度 & ACK 或规定数量 bit 已转移 & 模式、端序或帧边界错误 \\\tablerowrule
设备 FIFO & 读写指针、水位、溢出标志 & 样本被控制器取走 & 生产速度长期高于消费速度 \\\tablerowrule
DMA 环 & 描述符、所有权、时间戳范围 & completion 归还整批缓冲 & CPU 过早复用或 Cache 不可见 \\\tablerowrule
应用队列 & 单位、校准版本、序号与时间基 & 消费者确认处理或显式丢弃 & 把旧校准用于新量程 \\
\bottomrule
\end{longtable}

\section{十二、过载、丢样与重新同步}

缓冲区只能吸收短时速率差，不能修复长期过载。采样器以 10 ksample/s 产生 16-bit 三轴数据，未计头部时原始速率为 $10{,}000\times6=60$ kB/s；若控制器最长可能连续 40 ms 得不到总线服务，仅保存这段突发就需要至少 2400 B，再加水位中断延迟和 DMA 描述符切换裕量。队列满时必须定义策略：停止采样、覆盖最旧样本、丢弃最新样本或降低采样率，并通过序号间隙和溢出位让软件知道数据不连续。

错误恢复还要重新建立帧和时间关系。I2C 控制器可能因目标在一字节中途复位而看到总线被拉低；恢复流程可以产生规定数量时钟、发送 STOP 并重新初始化设备。SPI 没有地址和 ACK，片选异常或丢失一个时钟后，主从可能对字边界理解不同，通常要拉高片选结束帧并从已知命令重新开始。设备复位后，旧 FIFO、量程和校准状态都不能沿用；驱动应递增 generation，丢弃无法确定配置归属的残留样本。

从传感器到进程的完整 trace 因此是：物理量在采样沿被锁定，ADC 形成带状态的数字码，FIFO 保存顺序，接口控制器按帧取走，DMA 写入有所有权标记的内存，completion/中断通知驱动，驱动应用单位与校准并发布给进程。任何一层只要丢失序号、时间基或配置代际，数值即使“看起来合理”也可能不再代表正确的物理事件。




























\section{十三、Wi-Fi 与蓝牙：从内存帧到无线电波}

有线以太网把 MAC 帧交给 PCS、PMA 和 PMD；无线接口还要在 MAC 与天线之间完成信道接入、基带调制、数模转换、射频变频和功率放大。CPU 看到的仍是描述符、DMA 和完成项，但“介质正在被谁使用”不再由独占线对决定，而由共享频谱上的侦听、退避、确认和重传共同决定。

\begin{pdfdiagram}[x=1cm,y=1cm]
  \node[hw memory, minimum width=2.3cm] (mem) at (1.3,3.7) {主存\\帧/描述符};
  \node[hw control, minimum width=2.3cm] (mac) at (4.1,3.7) {MAC\\排队/退避/重传};
  \node[hw compute, minimum width=2.3cm] (bb) at (6.9,3.7) {基带\\编码/调制/OFDM};
  \node[hw compute, minimum width=2.3cm] (rf) at (9.7,3.7) {RF 收发机\\DAC/ADC/变频};
  \node[hw block, minimum width=2.3cm] (ant) at (12.5,3.7) {前端与天线\\PA/LNA/滤波};
  \draw[hw bus] (mem) -- (mac);
  \draw[hw bus] (mac) -- (bb);
  \draw[hw bus] (bb) -- (rf);
  \draw[hw bus] (rf) -- (ant);

  \node[hw state, minimum width=2.5cm] (air) at (12.5,1.0) {共享无线信道\\噪声/干扰/多径};
  \draw[hw bus] (ant) -- (air);
  \node[hw control, minimum width=2.8cm] (ctrl) at (8.3,1.0) {信道估计与速率控制\\RSSI/SNR/EVM/重传率};
  \draw[hw return] (air) -- (ctrl);
  \draw[hw return] (ctrl) -- (mac);
\end{pdfdiagram}
\diagramnote{无线发送与反馈是一个闭环。上排把内存帧逐层转换为射频能量，下排把信道测量、确认与重传结果送回速率控制和 MAC。吞吐下降可能来自队列、争用、基带误码、射频功率或天线环境，必须按层读取证据。}

\topic{发送路径把 bit 映射到频率、相位和幅度}

驱动写好发送描述符后，设备 DMA 读取帧和控制信息。MAC 添加序号、地址、校验和无线控制字段，并等待信道接入。基带先执行扰码、前向纠错编码和交织，再把 bit 组映射成调制符号；OFDM 把符号分配到许多正交子载波，经 IFFT 形成时域采样并加入循环前缀。DAC 把数字采样变成模拟基带或中频信号，混频器上变频到射频，功率放大器经滤波与天线把能量送入空间。

高阶调制让每个符号携带更多 bit，却要求更高信噪比和更低误差向量幅度。速率控制根据接收信号强度、信噪比、重传率和历史成功概率选择调制编码方案、信道宽度与空间流数量。最大标称速率只在足够信道质量、足够空间流和较少竞争时成立。

\topic{接收路径先恢复模拟波形，再判断 MAC 帧}

天线接收的信号先经带通滤波和低噪声放大器，再由混频器下变频；ADC 产生数字采样，自动增益控制避免过小或饱和。基带利用前导序列完成包检测、频偏估计、定时同步和信道估计，再执行 FFT、均衡、解调、解交织和纠错译码。只有通过物理层与 MAC 校验的帧，才被写入接收缓冲并产生完成项。

接收描述符耗尽时，即使射频和基带已正确解出帧，设备也没有合法 DMA 目标；相反，FCS、PHY 或重试计数上升时，增加软件接收环深度并不能修复信道问题。定位必须同时读取 RX ring 水位、MAC 丢弃、PHY 错误、RSSI/SNR、调制方案和重传计数。

\topic{关联与密钥建立发生在普通数据传输之前}

Wi-Fi 站点先扫描信道并发现网络，选择接入点后执行认证、关联和密钥协商，再安装用于数据帧保护的密钥。关联成功只证明控制状态已经建立；DHCP、路由和应用连接仍属于后续软件路径。漫游时，新接入点、密钥和队列代际必须在旧链路失效前后明确切换，避免旧重传帧被新连接错误接受。

蓝牙使用相同的“基带 + RF + 天线”物理骨架，但面向低功耗短包和跳频信道。控制器按连接事件唤醒，在指定频点交换少量包，然后重新休眠。低功耗来自缩短射频开启时间，不是让无线状态完全消失；时钟偏差、连接间隔、监督超时和重传仍决定连接能否持续。

\begin{longtable}{L{0.18\linewidth}L{0.31\linewidth}L{0.39\linewidth}}
\toprule
证据 & 更可能对应的边界 & 不应直接采取的动作 \\
\midrule
TX ring 长期积压 & MAC 争用、固件停顿或链路未关联 & 直接提高发射功率 \\\tablerowrule
重传率与 FCS 错误同时升高 & 干扰、多径、弱信号或频偏 & 只增加驱动队列深度 \\\tablerowrule
RSSI 高但吞吐低 & 同信道竞争、低 SNR、接收饱和或队列限制 & 仅依据信号格数判定链路良好 \\\tablerowrule
PHY 成功但主机无包 & RX 描述符、DMA、IOMMU 或中断路径 & 重新调天线匹配 \\\tablerowrule
关联反复重置 & 密钥、漫游、固件状态或时钟/电源问题 & 把每次断开都解释为射频弱 \\
\bottomrule
\end{longtable}

\section{十四、RDMA、RoCE 与 SmartNIC：网络设备直接访问应用缓冲区}

普通网络接收通常经过内核协议栈和 socket 缓冲，再复制或映射给应用。RDMA 把应用注册的内存区域、队列对和完成队列直接落实为 NIC 可访问的硬件对象，使远端请求可以把数据放入目标缓冲区，CPU 只处理建立、权限和完成。减少 CPU 参与并没有减少状态；它把更多所有权、顺序和错误恢复状态移动到了 RNIC。

\begin{pdfdiagram}[x=1cm,y=1cm]
  \node[hw state, minimum width=2.4cm] (app) at (1.3,3.4) {应用\\工作请求 WR};
  \node[hw memory, minimum width=2.4cm] (qp) at (4.2,3.4) {发送/接收队列\\QP 与 WQE};
  \node[hw control, minimum width=2.7cm] (rnic) at (7.3,3.4) {RNIC/SmartNIC\\传输、重传、权限};
  \node[hw block, minimum width=2.5cm] (wire) at (10.5,3.4) {Ethernet/RoCE\\或 InfiniBand};
  \draw[hw bus] (app) -- (qp);
  \draw[hw bus] (qp) -- (rnic);
  \draw[hw bus] (rnic) -- (wire);

  \node[hw memory, minimum width=2.7cm] (remote) at (10.5,1.0) {远端已注册内存\\rkey + 地址范围};
  \node[hw control, minimum width=2.7cm] (rrnic) at (7.3,1.0) {远端 RNIC\\IOMMU/页固定};
  \node[hw state, minimum width=2.4cm] (cq) at (4.2,1.0) {完成队列 CQ\\状态与代际};
  \node[hw state, minimum width=2.4cm] (done) at (1.3,1.0) {应用观察完成};
  \draw[hw bus] (wire) -- (remote);
  \draw[hw return] (remote) -- (rrnic);
  \draw[hw return] (rrnic) -- (cq);
  \draw[hw return] (cq) -- (done);
\end{pdfdiagram}
\diagramnote{一次 RDMA 操作的请求与完成环。上排由应用发布工作请求并经 RNIC 发送，下排由远端 RNIC 验证 rkey 和地址范围、访问已注册内存，再把完成返回本地 CQ。所谓“零拷贝”只减少 CPU 数据搬运，页固定、权限、队列槽、重传和完成代际仍必须闭合。}

\topic{内存注册把虚拟缓冲区变成设备可验证的能力}

应用注册内存时，内核固定或按需管理相关页，建立 IOMMU 映射，并向 RNIC 写入地址范围、权限和密钥。设备访问请求同时携带远端虚拟地址与 rkey；RNIC 查内存保护表，只有范围与权限都匹配才发起 DMA。rkey 不是加密密钥，而是不可猜测的能力标识；解除注册时必须阻止新请求、耗尽在途访问、撤销设备缓存的翻译和密钥，最后才能复用页框。

\topic{队列对与完成队列保存端到端顺序}

应用把 SEND、RECV、RDMA READ、RDMA WRITE 或原子操作写成 WQE，发布屏障后更新生产者指针。RNIC 读取 WQE、分段成网络包并维护包序、确认和重传。完成队列项表示该工作请求达到协议规定的本地完成边界，但不同操作的远端可见语义不同：本地 RDMA WRITE 完成通常说明数据已由传输层接受并抵达目标 RNIC/内存路径，不自动表示远端应用已经处理；需要应用级确认时仍要发送独立消息或更新双方约定的同步位置。

RoCE 把 RDMA 传输映射到以太网/IP 网络。无损网络可使用优先级流控与拥塞通知，但暂停帧会把一个拥塞点扩散到其他流量，拥塞控制仍要限制注入速率。丢包、ECN、重传超时、QP 状态和交换机队列应进入同一时间线，否则容易把网络拥塞误判为远端 CPU 延迟。

\topic{SmartNIC 与 DPU 把更多控制放到网卡侧}

SmartNIC 可以在网卡上执行虚拟交换、隧道封装、校验、加密和存储协议处理；DPU 进一步集成通用或专用处理核、内存与管理面。卸载后的数据路径可以绕过主机内核，但控制面仍要分配队列、IOMMU 域、密钥、带宽和复位代际。网卡固件异常时，主机必须能够隔离 DMA、收集设备日志并执行 FLR 或管理控制器复位。

\topic{PTP 硬件时间戳把“事件发生”固定在 MAC/PHY 边界}

软件在发送系统调用前后读取时钟，测到的是调度、队列和协议栈混合延迟。PTP 硬件时间戳在帧真正进入或离开 MAC/PHY 的固定位置锁存设备时钟，把可变排队部分排除。驱动把时间戳与描述符代际匹配后交给时间同步程序，程序估计时钟偏移与链路时延，再通过 PHC 或系统时钟调节接口校正频率和相位。时间戳若与错误队列项配对，偏差会表现为偶发的大幅跳变。

\section{十五、CAN 与 LIN：共享线上的确定性控制与错误恢复}

CAN 和 LIN 面向车辆与工业控制。它们的数据率低于 PCIe 和 Ethernet，却把电气鲁棒性、优先级仲裁、错误隔离和最坏响应时间放在首位。理解这两类总线可以看到：协议带宽较低并不意味着控制状态较少，反而常要求更明确的故障升级规则。

\begin{pdfdiagram}[x=1cm,y=1cm]
  \node[hw state, minimum width=2.4cm] (n0) at (1.6,4.1) {节点 A\\ID 0x120};
  \node[hw state, minimum width=2.4cm] (n1) at (4.7,4.1) {节点 B\\ID 0x080};
  \node[hw state, minimum width=2.4cm] (n2) at (7.8,4.1) {节点 C\\ID 0x300};
  \draw[BookMuted, line width=1.1pt] (1.6,2.65) -- (9.3,2.65);
  \draw[hw bus] (n0.south) -- (1.6,2.65);
  \draw[hw bus] (n1.south) -- (4.7,2.65);
  \draw[hw bus] (n2.south) -- (7.8,2.65);
  \node[hw block, minimum width=2.8cm] (bus) at (11.2,2.65) {CAN 差分总线\\显性 0 / 隐性 1};
  \draw[hw bus] (9.3,2.65) -- (bus);

  \node[hw control, minimum width=3.0cm] (arb) at (2.7,0.3) {逐 bit 仲裁\\发送 1 却读到 0 则退出};
  \node[hw control, minimum width=3.0cm] (err) at (6.8,0.3) {错误计数\\active/passive/bus-off};
  \node[hw state, minimum width=2.8cm] (winner) at (11.2,0.3) {最低 ID 继续发送\\其余节点等待重试};
  \draw[hw return] (bus.south) -- (11.2,1.45) -- (2.7,1.45) -- (arb.north);
  \draw[hw bus] (arb) -- (err);
  \draw[hw bus] (err) -- (winner);
\end{pdfdiagram}
\diagramnote{CAN 仲裁与错误隔离共享同一物理反馈：节点边发送边读取总线。仲裁阶段，读回更强的显性 0 表示自己失去优先级但不算错误；数据阶段，发送值与读回值不符会进入错误计数。协议根据当前阶段解释同一电气事实。}

\topic{CAN 用线与关系完成无破坏仲裁}

CAN 收发器把逻辑状态转换为 CAN\_H/CAN\_L 差分电压。显性 bit 0 可以覆盖隐性 bit 1，多个节点同时发送时，总线呈现逐 bit 线与结果。标识符从高位开始发送；某节点发送隐性 1 却读到显性 0，就知道另一个更低标识符帧拥有更高优先级，于是停止发送而不破坏获胜帧。最坏等待时间由高优先级帧频率、帧长度和总线利用率决定，不能只用平均带宽判断实时性。

数据阶段仍执行逐 bit 回读、位填充、CRC、格式和确认检查。节点发现错误后发送错误标志，使所有参与者丢弃该帧并重新发送。发送与接收错误计数器决定节点处于 error-active、error-passive 或 bus-off；故障节点若持续干扰总线，最终会被协议隔离。恢复 bus-off 必须满足平台规定的静默时间与管理策略，不能在电气短路仍存在时立即循环重启。

\topic{LIN 用主节点调度表换取更低成本}

LIN 通常使用单线加地、较低数据率和主从调度。主节点按调度表发送 break、同步字节和标识符，被指定的发布者随后发送数据与校验。没有逐 bit 多主仲裁，时间确定性来自主节点预先安排的时隙。节点可使用低成本 RC 时钟，通过同步字节校准当前帧的 bit 时间。

LIN 的低成本也带来边界：主节点故障会影响整条调度，单线对短路和地偏移更敏感，错误检测能力弱于 CAN。车辆常用 CAN 承担制动、动力等高可靠网络，用 LIN 承担车窗、座椅和传感器等低速支线；网关在两者之间转换信号时，必须定义更新时间、超时默认值和失联后的安全状态。

\section*{章末练习}

\begin{chapterexercise}{综合练习：1 至 2}
\begin{exercisesubproblem}{（a）1}
枚举
\exercisecheck{题目}{按顺序说明 USB 设备从连接到类驱动可提交传输之间的主要步骤。}
\end{exercisesubproblem}
\begin{exercisesubproblem}{（b）2}
带宽
\exercisecheck{题目}{为什么两个独立 endpoint 仍可能互相影响？列出至少两个共享资源。}
\end{exercisesubproblem}
\end{chapterexercise}

\begin{chapterexercise}{综合练习：3 至 4}
\begin{exercisesubproblem}{（a）3}
队列
\exercisecheck{题目}{xHCI 驱动写 TRB 后为何不能省略 DMA 同步与 doorbell 前的顺序约束？}
\end{exercisesubproblem}
\begin{exercisesubproblem}{（b）4}
复位
\exercisecheck{题目}{endpoint 复位时，怎样防止迟到事件或 DMA 访问已经复用的缓冲区？}
\end{exercisesubproblem}
\end{chapterexercise}

\begin{chapterexercise}{综合练习：5 至 6}
\begin{exercisesubproblem}{（a）5}
HID
\exercisecheck{题目}{解释 USB interrupt transfer 与 CPU 物理中断线的区别；按键释放报告丢失会出现什么现象？}
\end{exercisesubproblem}
\begin{exercisesubproblem}{（b）6}
音频
\exercisecheck{题目}{主机与设备标称采样率相同，为什么仍会逐渐欠载或溢出？反馈 endpoint 如何帮助稳定水位？}
\end{exercisesubproblem}
\end{chapterexercise}

\begin{chapterexercise}{综合练习：量化计算至采样边界}
\begin{exercisesubproblem}{（a）量化计算}
计算 12-bit、3.3 V ADC 的理想最低有效位
\exercisecheck{检查要点}{分辨率不等于精度}
\end{exercisesubproblem}
\begin{exercisesubproblem}{（b）采样边界}
说明抗混叠必须在 ADC 之前完成的原因
\exercisecheck{检查要点}{混叠后频率身份不可恢复}
\end{exercisesubproblem}
\end{chapterexercise}

\begin{chapterexercise}{时间戳}
比较采样沿时间和 CPU 中断时间
\exercisecheck{检查要点}{后者包含总线与调度抖动}
\end{chapterexercise}

\begin{chapterexercise}{缓冲预算}
计算正文三轴样本在 40 ms 停顿中的最低字节数
\exercisecheck{检查要点}{还要为控制延迟留裕量}
\end{chapterexercise}

\begin{chapterexercise}{接口恢复}
分别给出 I2C 总线卡住和 SPI 帧失步的重新同步入口
\exercisecheck{检查要点}{从已知帧边界重新开始}
\end{chapterexercise}

\begin{chapterexercise}{代际隔离}
说明传感器复位后为何要丢弃旧 FIFO 样本
\exercisecheck{检查要点}{配置与校准归属可能已改变}
\end{chapterexercise}

\begin{chapterexercise}{UART 帧分析}
给定 RX 波形和标称波特率，数出起始位、数据位、停止位，还原字节并核对波特率误差
\exercisecheck{预期结果}{帧结构和误差判断都有依据}
\end{chapterexercise}

\begin{chapterexercise}{I2C 事务序列}
写出主机从 7 位地址 \code{0x50} 的从机寄存器 \code{0x10} 读 2 字节的完整序列：START、地址字节、读写位、ACK/NACK、repeated START、STOP
\exercisecheck{预期结果}{\code{0xA0}/\code{0xA1} 正确，最后一字节 NACK，repeated START 位置正确}
\end{chapterexercise}

\section*{参考解答与要点}

\begin{chapteranswer}{综合练习：1 至 2}
\begin{answersubproblem}{（a）1}
枚举
\answercheck{解答要点}{检测连接与供电，端口复位，在默认地址读取基本描述符，分配地址，读取配置/接口/endpoint 描述符，选择配置和 alternate setting，绑定类驱动。}
\end{answersubproblem}
\begin{answersubproblem}{（b）2}
带宽
\answercheck{解答要点}{endpoint 可能共享设备上行、hub 上行、周期调度预算、主控制器 ring、DMA 和中断处理，因此逻辑独立不代表性能与故障域独立。}
\end{answersubproblem}
\end{chapteranswer}

\begin{chapteranswer}{综合练习：3 至 4}
\begin{answersubproblem}{（a）3}
队列
\answercheck{解答要点}{必须先让 TRB 和缓冲内容对控制器可见，再通知控制器取用；否则控制器可能看到 doorbell 却读到旧描述符或旧数据。}
\end{answersubproblem}
\begin{answersubproblem}{（b）4}
复位
\answercheck{解答要点}{先停止 endpoint 并等待控制器确认，取消或隔离旧 TRB，改变 cycle/generation，处理 event ring 中旧事件，确认无 DMA 后才释放和复用缓冲。}
\end{answersubproblem}
\end{chapteranswer}

\begin{chapteranswer}{综合练习：5 至 6}
\begin{answersubproblem}{（a）5}
HID
\answercheck{解答要点}{interrupt endpoint 由主机周期轮询，CPU 中断是主控制器报告完成的机制。释放报告丢失可能形成“粘键”，需后续完整状态或重新同步纠正。}
\end{answersubproblem}
\begin{answersubproblem}{（b）6}
音频
\answercheck{解答要点}{独立晶振存在频差；反馈提供实际消费速率，主机按服务间隔微调样本数，使长期平均生产率匹配并把水位保持在目标附近。}
\end{answersubproblem}
\end{chapteranswer}

\begin{chapteranswer}{综合练习：量化计算至采样边界}
\begin{answersubproblem}{（a）量化计算}
$3.3/4096\approx0.806$ mV/LSB；实际精度还受参考、噪声、偏移与非线性影响。
\answercheck{必须保留的检查点}{不把码宽直接当准确度}
\end{answersubproblem}
\begin{answersubproblem}{（b）采样边界}
高频分量采样后会折叠到基带，与真实低频产生相同样本序列；数字处理无法区分二者。
\answercheck{必须保留的检查点}{模拟滤波位于采样之前}
\end{answersubproblem}
\end{chapteranswer}

\begin{chapteranswer}{时间戳}
采样沿代表物理量被锁定的时刻；中断时间还包含转换、FIFO、水位、总线、合并和调度延迟。
\answercheck{必须保留的检查点}{时间戳要注明时钟域}
\end{chapteranswer}

\begin{chapteranswer}{缓冲预算}
每样本 6 B，10 ksample/s 在 0.04 s 内产生 $6\times10{,}000\times0.04=2400$ B。
\answercheck{必须保留的检查点}{这是未计裕量的下限}
\end{chapteranswer}

\begin{chapteranswer}{接口恢复}
I2C 尝试释放时钟并发 STOP 后重配；SPI 拉高片选结束残帧，再从已知命令重新开始。
\answercheck{必须保留的检查点}{不能从半帧继续解释}
\end{chapteranswer}

\begin{chapteranswer}{代际隔离}
复位可能恢复默认量程、采样率和校准，残留码无法确定属于哪套配置；generation 变化后应拒绝旧样本。
\answercheck{必须保留的检查点}{数值合理不代表语义正确}
\end{chapteranswer}

\begin{chapteranswer}{UART 帧分析}
先找起始位下降沿，再按位时间中点逐位采样：1 起始位、8 数据位（低位在前）、1 停止位；还原字节后核对实测位时间与标称波特率的偏差，双方每一侧误差一般不超过约 2\%，收发合计留有约 5\% 的理论上限。
\answercheck{必须保留的要点}{位边界、采样点和误差都要核对。}
\end{chapteranswer}

\begin{chapteranswer}{I2C 事务序列}
START $\rightarrow$ \code{0xA0}（\code{0x50} 左移一位、读写位 0 表示写）$\rightarrow$ 从机 ACK $\rightarrow$ \code{0x10}（寄存器号）$\rightarrow$ ACK $\rightarrow$ repeated START $\rightarrow$ \code{0xA1}（读写位 1 表示读）$\rightarrow$ ACK $\rightarrow$ 读第 1 字节 $\rightarrow$ 主机 ACK $\rightarrow$ 读第 2 字节 $\rightarrow$ 主机 NACK $\rightarrow$ STOP。
\answercheck{必须保留的要点}{先写寄存器地址，repeated START 转读；最后字节 NACK 通知从机结束。}
\end{chapteranswer}
